% =====================================================================
%  Capitolo 16 — I salvataggi: Save Game, Game Instance e le monete
%  che non si dimenticano, più una mattinata di feature da team (la
%  quarta lezione di programmazione con i Blueprint di Unreal Engine)
%  (Game Design, primo anno)
%  Fonte: trascrizione della lezione GD16 (Game Design, primo anno),
%  quarta lezione di programmazione, diretta prosecuzione dei capitoli
%  7, 12 e 13 (stesso docente). È la lezione "del 23" annunciata nella
%  chiusura del capitolo 13, l'ultima prima della consegna del
%  platformer. Lezione ibrida: una parte della classe in aula, una parte
%  collegata da remoto (il docente, in chiusura, si rivolge a "quelli
%  online").
%  Contenuti, nell'ordine della lezione: la chiacchierata d'apertura
%  (Unreal Engine 6.0 e il ritmo delle versioni, la build per browser
%  tolta e rimessa e perché non può essere un plugin, l'engine scritto
%  da un'IA in un giorno e mezzo, la storia degli engine: dagli engine
%  fatti in casa a Unity, Unreal e CryEngine, e forse ritorno; Godot e il
%  3D); il punto sui task e su GitHub; il template in prima persona e
%  come installare i template dal launcher (template e pacchetti di
%  funzionalità, risparmiare spazio); i salvataggi: perché non esiste un
%  "tasto salva tutto"; la classe Save Game (PlayerSaves, una o due per
%  gioco, la variabile Score); la Game Instance rivista (che cosa ci va
%  oltre ai salvataggi: achievement, musica tra un livello e l'altro;
%  perché sopravvive ai cambi di mappa; Event Init); i nodi del
%  salvataggio (crea, controlla, carica, cancella, salva) e le varianti
%  asincrone (nodi istantanei e nodi latenti, l'orologino, l'autosave
%  degli open world, l'icona che lampeggia); Does Save Game Exist, Branch,
%  Create Save Game Object e Load Game from Slot; lo slot name e la
%  variabile stringa; la reference al salvataggio e le variabili che
%  hanno per tipo un oggetto; il cast al tipo di salvataggio; Maps &
%  Modes; Save Game to Slot e il custom event SaveGame; le monete (il
%  Blueprint Coins, il cilindro, l'overlap, il cast al personaggio, il
%  custom event CoinGrabbed nel player, Destroy Actor); la "torre di
%  nodi": Get Game Instance, cast, Get Save, Get/Set Score, il cast puro;
%  la prova (1-2-3 che si ripete senza salvare, poi 4-5-6, 7-8-9 anche
%  dopo aver chiuso); il file .sav in Saved/SaveGames; il riepilogo; le
%  domande (dove si trova il SaveGame, il print continuo con l'Event Tick,
%  quali nodi diventano puri); l'anticipazione del best score nel
%  prossimo task; l'esercitazione "da azienda" con nove feature da
%  dividersi (la scala con la distanza, il jetpack, la torcia con la
%  batteria, il ritorno al punto sicuro, l'oggetto che c'è il 35% delle
%  volte, la torretta, i tre spawn con l'attesa di 5 secondi, la zona
%  da 10 HP al secondo, il modello che cambia ogni secondo); i
%  suggerimenti durante il lavoro (Get Player Pawn, la doppia mesh del
%  personaggio in prima persona, l'input solo sui Pawn posseduti, a
%  differenza di Unity); la chiusura.
%  Integrazioni dai libri "Blueprints Visual Scripting for Unreal Engine
%  5" (Romero, Sewell) e "The Art of Game Design: A Book of Lenses"
%  (Jesse Schell, terza edizione), segnalate come approfondimenti.
%
%  Pacchetti richiesti nel preambolo del libro (già presenti nel
%  main.tex): \usepackage{booktabs}, \usepackage{amsmath},
%  \usepackage{enumitem}, \usepackage{tikz} con
%    \usetikzlibrary{shapes.geometric,shapes.symbols,shapes.misc,
%      arrows.meta,positioning,calc,fit}
%  (le librerie vanno caricate nel preambolo, non nel capitolo).
%  Tutte le etichette di questo capitolo hanno il prefisso "gd16:"
%  (gd16:sec:..., gd16:fig:..., gd16:tab:...) e tutti gli stili TikZ
%  "propri" il prefisso "gd16", per non collidere con gli altri
%  capitoli. I grafi Blueprint usano le macro di bp-tikz.tex (la piccola
%  libreria in tasks/gd7-blueprint/), incluse qui sotto così che il
%  capitolo sia autosufficiente; sono protette da un controllo, quindi
%  se un altro capitolo le ha già definite non vengono ridefinite.
%
%  Nota sulla fonte: la trascrizione automatica è buona nella
%  chiacchierata iniziale e nella spiegazione dei salvataggi (circa un
%  terzo della registrazione), ma tutta la seconda parte della lezione,
%  quella del lavoro individuale sulle feature, è quasi interamente
%  rovinata: il riconoscimento vocale ha riempito i silenzi con frasi in
%  loop ("Vediamo cosa succede se qualcun altro non ha problemi",
%  "non voglio fare come una collezione", "la struttura del giardino",
%  "dopo tutto, forse adesso non sono le stesse lerecchie") e con
%  lunghissimi brani captati da una fonte esterna o inventati (frasi
%  sulla "dottrina sul mercato del lavoro", sulla "Corte Suprema di
%  Internet", sulla "città-stato", sui "funzionari della società", "la
%  migliore dal 2005", "il numero di persone è molto più basso dopo il
%  raggio di persone", "gli Stati Uniti", "gli egiziani", "la signora
%  Nicolai", "Vittorio Prodi", "i patrioti della società civile", "il
%  paziente"): sono artefatti, omessi. Di quella parte sopravvivono solo
%  tre suggerimenti dati ad alta voce alla classe e la frase di
%  chiusura, che la trascrizione taglia a metà. Diversi nomi sono
%  storpiati ("Arrio", "Ariadne", "Real" per Unreal; "BravaLoser" per
%  browser; "Cray Engine" per CryEngine; "Bluetooth class" per Blueprint
%  Class; "gate distance", "distance" per Game Instance; "Scorra" per
%  Score; "player sense" per PlayerSaves; "lot save, load game from
%  flot" per Load Game from Slot; "Syncseg" per il salvataggio
%  asincrono; "Boxer Game" per il custom event SaveGame; "a sette" per
%  set; "destroy hardware" per Destroy Actor; "eventic" per Event Tick;
%  "scamotaggio" per escamotage; "Tobiquelli" per "per quelli"): dove il
%  riferimento è ricostruibile con certezza lo si riporta nella forma
%  corretta; dove non lo è, lo si segnala in nota invece di indovinare.
%  Le figure sono ricostruzioni schematiche dei grafi costruiti a
%  lezione e dei concetti spiegati, non screenshot.
% =====================================================================

% --- Macro per i grafi Blueprint (da tasks/gd7-blueprint/bp-tikz.tex) --
% Richiedono TikZ e la libreria calc, già caricati nel preambolo.
\ifcsname bpnode\endcsname\else
\pgfdeclarelayer{bpbg}
\pgfdeclarelayer{bpcomment}
\pgfdeclarelayer{bpwire}
\pgfsetlayers{bpbg,bpcomment,bpwire,main}

% --- colori dell'editor ----------------------------------------------
\definecolor{bpBg}{HTML}{262626}
\definecolor{bpGridFine}{HTML}{2E2E2E}
\definecolor{bpGridBold}{HTML}{1A1A1A}
\definecolor{bpBody}{HTML}{141414}
\definecolor{bpOutline}{HTML}{060606}
\definecolor{bpText}{HTML}{F0F0F0}
\definecolor{bpSub}{HTML}{B4B4B4}
\definecolor{bpValBg}{HTML}{0C0C0C}
\definecolor{bpValBorder}{HTML}{5A5A5A}

% header dei nodi
\definecolor{bpHEvent}{HTML}{8E1B1B}   % eventi (rosso)
\definecolor{bpHFunc}{HTML}{1F5A8C}    % funzioni con esecuzione (blu)
\definecolor{bpHPure}{HTML}{3F7A2E}    % funzioni pure, senza esecuzione (verde)
\definecolor{bpHCast}{HTML}{14706F}    % cast (verde acqua)
\definecolor{bpHFlow}{HTML}{5C5C5C}    % controllo del flusso, macro (grigio)
\definecolor{bpHVar}{HTML}{3A3A3A}     % SET di una variabile

% pin e fili, un colore per tipo di dato (come nell'editor)
\definecolor{bpPinexec}{HTML}{FFFFFF}
\definecolor{bpPinbool}{HTML}{A00000}
\definecolor{bpPinint}{HTML}{1FE3AF}
\definecolor{bpPinfloat}{HTML}{9DF53D}
\definecolor{bpPinstr}{HTML}{F72BD3}
\definecolor{bpPinname}{HTML}{C98AF7}
\definecolor{bpPintext}{HTML}{E2779A}
\definecolor{bpPinvec}{HTML}{F8C624}
\definecolor{bpPinrot}{HTML}{A0B3FC}
\definecolor{bpPintrans}{HTML}{F77A10}
\definecolor{bpPinobj}{HTML}{19A7F2}
\definecolor{bpPinclass}{HTML}{7B2BC4}
\definecolor{bpPinsoft}{HTML}{98F2FF}
\definecolor{bpPinstruct}{HTML}{1458C8}
\definecolor{bpPinwild}{HTML}{8A8A8A}
\definecolor{bpPinenum}{HTML}{0D7A6B}
\definecolor{bpPinskip}{HTML}{141414}

\def\bpRow{0.42}
\def\bpExecStr{exec}
\def\bpSkipStr{skip}
\def\bpCompactStr{compact}

\tikzset{
  bpTitle/.style={text=bpText, font=\sffamily\scriptsize\bfseries,
                  inner sep=0pt, align=left},
  bpTitleC/.style={text=bpSub, font=\sffamily\large\bfseries, inner sep=0pt},
  bpLabel/.style={text=bpText, font=\sffamily\tiny, inner sep=0pt},
  bpNoteSty/.style={text=bpSub, font=\sffamily\tiny\itshape, align=left,
                    inner sep=1pt},
}

\newcommand{\bpval}[1]{{\setlength{\fboxsep}{0.8pt}\setlength{\fboxrule}{0.3pt}%
  \fcolorbox{bpValBorder}{bpValBg}{\textcolor{bpText}{\,#1\,}}}}
\newcommand{\bpsub}[1]{{\color{bpSub}\tiny\itshape\mdseries #1}}

% conta gli elementi non vuoti di una lista "tipo/etichetta, ..."
\newcommand{\bpcountlist}[1]{%
  \gdef\bpres{0}%
  \foreach \bpt/\bpl in {#1}{%
    \ifx\bpt\empty\else
      \pgfmathtruncatemacro\bptmp{\bpres+1}\global\let\bpres\bptmp
    \fi}}

% disegna un pin: tipo, x, y
\newcommand{\bpXpin}[3]{%
  \edef\bpPT{#1}%
  \ifx\bpPT\bpExecStr
    \fill[bpPinexec] (#2-0.085,#3+0.085) -- (#2+0.02,#3+0.085)
       -- (#2+0.105,#3) -- (#2+0.02,#3-0.085) -- (#2-0.085,#3-0.085) -- cycle;
  \else\ifx\bpPT\bpSkipStr
  \else
    \fill[bpPin#1] (#2,#3) circle (0.07);
    \draw[bpPin#1!60!black, line width=0.3pt] (#2,#3) circle (0.07);
  \fi\fi}

% il nodo generico
\newcommand{\bpnode}[8][1]{%
  \bpcountlist{#7}\edef\bpNin{\bpres}%
  \bpcountlist{#8}\edef\bpNout{\bpres}%
  \pgfmathsetmacro\bpRows{max(\bpNin,\bpNout,1)}%
  \edef\bpHC{#4}%
  \ifx\bpHC\bpCompactStr
    \def\bpHead{0.06}%
  \else
    \pgfmathsetmacro\bpHead{0.14+0.34*#1}%
  \fi
  \pgfmathsetmacro\bpH{\bpHead+\bpRows*\bpRow+0.12}%
  \begin{scope}[shift={(#3)}]
    \fill[black, opacity=0.45, rounded corners=4pt]
         (0.07,-0.07) rectangle ++(#6,-\bpH);
    \fill[bpBody, rounded corners=4pt] (0,0) rectangle (#6,-\bpH);
    \ifx\bpHC\bpCompactStr
      \node[bpTitleC] at ({#6-0.52},{-\bpH/2}) {#5};
    \else
      \begin{scope}
        \clip[rounded corners=4pt] (0,0) rectangle (#6,-\bpH);
        \shade[left color=#4, right color=#4!30!bpBody]
              (0,0) rectangle (#6,-\bpHead);
      \end{scope}
      \node[bpTitle, anchor=west] at (0.14,{-\bpHead/2}) {#5};
    \fi
    \draw[bpOutline, line width=0.6pt, rounded corners=4pt]
         (0,0) rectangle (#6,-\bpH);
    \coordinate (#2-c)  at ({#6/2},{-\bpH/2});
    \coordinate (#2-n)  at ({#6/2},0);
    \coordinate (#2-s)  at ({#6/2},-\bpH);
    \coordinate (#2-w)  at (0,{-\bpH/2});
    \coordinate (#2-e)  at (#6,{-\bpH/2});
    \coordinate (#2-sw) at (0,-\bpH);
    \coordinate (#2-ne) at (#6,0);
    \foreach \bpt/\bpl [count=\bpk] in {#7}{%
      \ifx\bpt\empty\else
        \pgfmathsetmacro\bpy{-\bpHead-(\bpk-0.5)*\bpRow-0.06}%
        \bpXpin{\bpt}{0.17}{\bpy}%
        \coordinate (#2-i\bpk) at (0.17,\bpy);
        \node[bpLabel, anchor=west] at (0.33,\bpy) {\bpl};
      \fi}%
    \foreach \bpt/\bpl [count=\bpk] in {#8}{%
      \ifx\bpt\empty\else
        \pgfmathsetmacro\bpy{-\bpHead-(\bpk-0.5)*\bpRow-0.06}%
        \bpXpin{\bpt}{#6-0.17}{\bpy}%
        \coordinate (#2-o\bpk) at (#6-0.17,\bpy);
        \node[bpLabel, anchor=east] at (#6-0.33,\bpy) {\bpl};
      \fi}%
  \end{scope}}

% il GET di una variabile: una "pillola" con un'uscita
\newcommand{\bpget}[5]{%
  \begin{scope}[shift={(#2)}]
    \fill[black, opacity=0.45, rounded corners=6pt]
         (0.06,-0.06) rectangle ++(#5,-0.46);
    \begin{scope}
      \clip[rounded corners=6pt] (0,0) rectangle (#5,-0.46);
      \shade[left color=bpPin#3!40!bpBody, right color=bpBody]
            (0,0) rectangle (#5,-0.46);
    \end{scope}
    \draw[bpOutline, line width=0.6pt, rounded corners=6pt]
         (0,0) rectangle (#5,-0.46);
    \node[bpTitle, anchor=west] at (0.14,-0.23) {#4};
    \bpXpin{#3}{#5-0.17}{-0.23}%
    \coordinate (#1-o1) at (#5-0.17,-0.23);
    \coordinate (#1-n) at ({#5/2},0);
    \coordinate (#1-s) at ({#5/2},-0.46);
  \end{scope}}

% un filo (sotto i nodi)
\newcommand{\bpwire}[3]{%
  \begin{pgfonlayer}{bpwire}
    \edef\bpWT{#1}%
    \ifx\bpWT\bpExecStr\def\bpWW{1.15pt}\else\def\bpWW{0.95pt}\fi
    \draw[bpPin#1, line width=\bpWW]
      let \p1=(#2), \p2=(#3),
          \n1={ifthenelse(\x2<\x1, 45pt, max(12pt, abs(\x2-\x1)*0.45))} in
      (#2) .. controls +(\n1,0) and +(-\n1,0) .. (#3);
  \end{pgfonlayer}}

% riquadro di commento
\newcommand{\bpcomment}[3]{%
  \begin{pgfonlayer}{bpcomment}
    \fill[white, opacity=0.07, rounded corners=3pt] (#1) rectangle (#2);
    \draw[white!45!bpBg, rounded corners=3pt, line width=0.4pt] (#1) rectangle (#2);
    \begin{scope}
      \clip[rounded corners=3pt] (#1) rectangle (#2);
      \path (#1) coordinate (bpcA) (#2) coordinate (bpcB);
      \path ($(bpcA)+(0,-0.42)$) coordinate (bpcH);
      \fill[white, opacity=0.10] (bpcA) rectangle (bpcB |- bpcH);
    \end{scope}
    \node[text=bpText, font=\sffamily\scriptsize\bfseries, anchor=north west,
          inner sep=3pt] at (#1) {#3};
  \end{pgfonlayer}}

\newcommand{\bpnote}[2]{\node[bpNoteSty, anchor=north west] at (#1) {#2};}

% l'ambiente grafo
\newenvironment{bpgraph}[1][]{%
  \begin{tikzpicture}[#1]}{%
  \begin{pgfonlayer}{bpbg}
    \path (current bounding box.south west) ++(-0.35,-0.35) coordinate (bpSW);
    \path (current bounding box.north east) ++(0.35,0.35) coordinate (bpNE);
    \fill[bpBg] (bpSW) rectangle (bpNE);
    \begin{scope}
      \clip (bpSW) rectangle (bpNE);
      \draw[bpGridFine, line width=0.3pt, step=0.25] (bpSW) grid (bpNE);
      \draw[bpGridBold, line width=0.5pt, step=2] (bpSW) grid (bpNE);
    \end{scope}
  \end{pgfonlayer}
  \end{tikzpicture}}
\fi

\chapter[I salvataggi: Save Game e Game Instance]{I salvataggi: Save Game,
Game Instance e le monete che non si dimenticano}
\label{gd16:cap:salvataggi}

Il capitolo~\ref{gd13:cap:git-componenti} si era chiuso con un appuntamento:
la lezione ``del 23'', l'ultima prima della consegna del platformer
(sezione~\ref{gd12:sec:task}). Eccola: è la \textbf{quarta lezione di
programmazione} con i Blueprint, con lo stesso docente dei
capitoli~\ref{gd7:cap:blueprint}, \ref{gd12:cap:array-timeline}
e~\ref{gd13:cap:git-componenti}. Nel capitolo precedente la Game Instance era
stata presentata con un avvertimento preciso: \emph{non è un salvataggio},
perché quando si chiude il gioco i suoi dati si perdono. Questa lezione chiude
il cerchio e mostra come si scrivono davvero dei dati \textbf{sul disco}, in
modo che sopravvivano alla chiusura del gioco e perfino allo spegnimento del
computer.

È anche la prima volta, lo dice il docente stesso, che si
\textbf{interconnettono così tanto i Blueprint}: una moneta che chiama il
personaggio, un personaggio che chiede alla Game Instance il salvataggio, un
salvataggio che contiene il punteggio, ``variabili con dentro variabili con
dentro variabili''. Il salvataggio delle monete costruito a lezione, avverte,
non è un esempio fine a sé stesso: ``questa cosa qui che viene salvato il
\textbf{best score}, ve lo anticipo già, ci sarà nel \textbf{prossimo task}''.

La seconda parte è pratica e cambia formato: il docente simula un
\textbf{team di sviluppo} in un'azienda, con un elenco di \textbf{nove
feature} da implementare e da dividersi, una a testa. La registrazione di
quella parte, in cui ciascuno lavora sul proprio progetto, è quasi tutta
illeggibile (si veda la nota in testa al file sorgente): ne restano l'elenco
delle feature e tre suggerimenti dati ad alta voce. Il capitolo li riporta e,
per ciascuna feature, propone in un riquadro di approfondimento una possibile
strada, chiaramente segnalata come integrazione.

La struttura della lezione, e quindi del capitolo, è questa:
\begin{enumerate}
  \item la chiacchierata d'apertura: Unreal Engine~6, la build per browser,
        l'engine scritto da un'IA e la storia degli engine
        (sezione~\ref{gd16:sec:apertura});
  \item il punto sui task e su GitHub, il \textbf{template in prima persona} e
        come si installano i template dal launcher
        (sezioni~\ref{gd16:sec:task-github} e~\ref{gd16:sec:template});
  \item perché non esiste un ``tasto salva tutto'' e la classe
        \textbf{Save Game} (sezioni~\ref{gd16:sec:perche}
        e~\ref{gd16:sec:savegame});
  \item la \textbf{Game Instance} rivista: che cosa ci va e perché
        (sezione~\ref{gd16:sec:game-instance});
  \item i nodi del salvataggio e i nodi \textbf{asincroni}
        (sezione~\ref{gd16:sec:nodi});
  \item il caricamento all'avvio: controlla, crea o carica, la reference e il
        cast (sezioni~\ref{gd16:sec:init} e~\ref{gd16:sec:reference});
  \item dirlo all'engine e \textbf{salvare davvero}: Maps \& Modes, Save Game
        to Slot e il custom event \emph{SaveGame}
        (sezione~\ref{gd16:sec:salvare});
  \item le \textbf{monete}: il Blueprint, il custom event nel personaggio e la
        ``torre di nodi'' (sezioni~\ref{gd16:sec:monete}
        e~\ref{gd16:sec:torre});
  \item la prova, il file \texttt{.sav} sul disco e il riepilogo
        (sezione~\ref{gd16:sec:prova});
  \item le domande della classe (sezione~\ref{gd16:sec:domande});
  \item l'esercitazione ``da azienda'' con le nove feature e i suggerimenti
        durante il lavoro (sezioni~\ref{gd16:sec:esercitazione}
        e~\ref{gd16:sec:suggerimenti});
  \item la chiusura (sezione~\ref{gd16:sec:chiusura}).
\end{enumerate}

% ---------------------------------------------------------------------
\section{In apertura: engine vecchi, engine nuovi}
\label{gd16:sec:apertura}

La registrazione parte a conversazione già avviata, mentre la classe si
sistema e il docente aspetta che si colleghi chi segue da remoto. Si parla
di engine, ed è una chiacchierata che vale la pena riportare perché tocca un
tema che riguarda da vicino chi farà questo mestiere nei prossimi anni.

\subsection{Unreal Engine 6 e il ritmo delle versioni}

Il primo argomento è l'uscita di \textbf{Unreal Engine~6.0}: ``in teoria''
esce \textbf{entro fine anno}. C'è già una \emph{preview}, che chi vuole
può scaricare, ma secondo quanto si dice in aula è stata fatta uscire ``di
fretta'', in concomitanza con un \emph{contest} per il quale si voleva che si
usasse la \textbf{5.9}.\footnote{Il passaggio è molto frammentato nella
trascrizione (``Rolevo la sei punto zero. Vero esce, in teoria, entro fine anno
così a Ber''): si riporta solo ciò che è ricostruibile, cioè che la 6.0 è
attesa per fine anno e che una preview è già disponibile. Non è chiaro dalla
registrazione se a parlare sia il docente o uno studente, né quale fosse il
contest.} Il ritmo, comunque, è noto: ``tendenzialmente c'è una versione
\textbf{ogni sei mesi}'', quindi la prossima ``dovrebbe uscire verso
gennaio, credo, più o meno''. (Nel capitolo~\ref{gd12:cap:array-timeline}
il docente aveva raccontato di stare già compilando Unreal Engine~6 dal
codice sorgente su GitHub per provarne le funzioni.)

\subsection{La build per browser, tolta e rimessa}

Si torna sulla novità della 5.9 già vista nel
capitolo~\ref{gd12:cap:array-timeline}: la \textbf{build per browser}, cioè
la possibilità di pubblicare un gioco giocabile direttamente nel
browser.\footnote{Nella trascrizione ``BravaLoser'': dal contesto (``aggiornare
l'engine per controllarla sempre anche su\dots'') si tratta quasi certamente
del browser.} La storia, ricostruita in aula, è questa:
\begin{itemize}
  \item la funzione era stata \textbf{tolta ``un sacco di anni fa''}, perché
        si pensava che ci fosse ``troppo poco mercato'' ed era ``un sacco di
        casino'' aggiornare l'engine controllando ogni volta che tutto
        funzionasse anche sul browser;
  \item la versione per browser che esisteva era ormai ``vecchissima'': Epic
        l'aveva rilasciata dicendo, in sostanza, ``continuatela a sviluppare
        voi come community''; ``nessuno l'aveva fatto, e quindi è morta lì'';
  \item adesso Epic l'ha rimessa, ``e quindi ci sta, è giusto''. ``Feature
        che per una volta ha anche senso.''
\end{itemize}

\paragraph{Perché non può essere un plugin.} Il docente aggiunge un
retroscena. Ne ha parlato con uno sviluppatore che crea \textbf{plugin} su
\textbf{Fab} (il marketplace visto nella sezione~\ref{gd12:sec:fab}), ``molto
bravo, mi vende anche tanta roba'', che gli ha spiegato che la build per
browser \textbf{non si può realizzare come plugin}, e quindi nemmeno vendere:
``per come è strutturato l'engine, non puoi creare un plugin che aggiunge la
build per il browser''. È ``una limitazione dell'engine proprio''. Da qui il
ragionamento: se fosse fattibile come plugin, ``lo farebbero tutti e lo
venderebbero''; non essendolo, nessuno lo faceva, e alla fine ha dovuto
rimetterla Epic.

\paragraph{La 5.9 e l'IA.} Sempre sulla 5.9 circola la voce che sia stata
fatta in gran parte con l'aiuto dell'intelligenza
artificiale.\footnote{La trascrizione dice ``l'hanno fatta tipo tutta col
diablo, così hanno detto'', seguito da un frammento incomprensibile
(``Finché mettono almeno 15 minuti, eh no, non è fatti''). Il riferimento
all'IA è un'interpretazione resa plausibile dal discorso che segue
immediatamente, ma non è verificabile dalla registrazione.} Il docente lo
riconosce come ``una digressione un pochettino'', ma la porta avanti con un
esperimento personale.

\subsection{L'engine scritto in un giorno e mezzo}

``Qualche settimana fa'' aveva a disposizione dei \textbf{crediti a tempo} su
un servizio di intelligenza artificiale\footnote{Nella trascrizione ``crediti
limitati su Cloud'': può trattarsi di un servizio cloud generico o del nome di
un assistente di IA storpiato dal riconoscimento vocale; dal racconto è chiaro
solo che si tratta di un agente di IA capace di scrivere codice in autonomia.},
``troppissimi'', e non sapeva che farne. Così ha provato:
\begin{quote}
Boh, facciamo così: fammi un altro engine, tipo Unreal. [\dots] Lo lascio
andare. Un giorno e mezzo dopo torno. Ha fatto veramente un gran bell'engine,
pazzesco, veramente un mostro.
\end{quote}
Il risultato, a sentire il docente, è notevole:
\begin{itemize}
  \item \textbf{graficamente} ``sembrava Unreal, tranquillamente'';
  \item aveva rimesso dentro \textbf{i Blueprint}, ``ma più ottimizzati,
        assurdo'';
  \item era \textbf{già pronto per il multiplayer};
  \item supportava il \textbf{2D};
  \item funzionava \textbf{già su browser}.
\end{itemize}

\subsection{Una breve storia degli engine}

Da qui il docente allarga il discorso: ``aspetta, prendiamola un attimino
sulla parte storica''.
\begin{enumerate}
  \item \textbf{All'inizio}, ogni casa di sviluppo si faceva \textbf{il proprio
        engine}: ``facevi l'engine per fare quel gioco'', non un engine per
        tutti i giochi.
  \item Poi, man mano, è diventato sempre più complesso \textbf{stare dietro ai
        requisiti}, e ogni volta bisognava \textbf{reimplementare} da capo il
        funzionamento delle varie parti\footnote{Nella trascrizione ``come
        funzionava il GACI'': il termine non è ricostruibile con certezza (può
        riferirsi a un sottosistema dell'engine, per esempio
        l'illuminazione globale, la \emph{GI}). Il senso del passaggio, cioè
        che ogni studio doveva reimplementare da zero sistemi complessi, è
        comunque chiaro.}: ``tutto quanto era un casino''. Ci si è quindi
        spostati sugli \textbf{engine comuni}: \textbf{Unity},
        \textbf{Unreal}, \textbf{CryEngine} e altri.
  \item \textbf{Oggi} questa scelta ``sta scricchiolando''. Se un'IA può
        scriverti l'engine, non c'è più bisogno di tenersi un Unreal ``che ha un
        sacco di \emph{bloat}'', cioè ``un sacco di roba che non ti serve'' e
        che ``ti appesantisce'': si può \textbf{tornare a farsi l'engine per il
        proprio gioco}.
\end{enumerate}
A favore di questo ritorno c'è anche un motivo economico: con Unity, Unreal e
simili, ``se cominci a fare un po' di vendita, devi anche pagarli'', mentre un
engine proprio ``è tutto quanto tuo, ti fai quel che vuoi''.

\paragraph{Godot e il 3D.} Per ora, tra gli engine gratuiti e
\emph{open source}, c'è \textbf{Godot} (citato già nel
capitolo~\ref{cap:introduzione-game-design}), ``che è carino, è carino, però con il 3D non
riesce veramente a fare niente''. Il giudizio del docente è netto: crede
sinceramente che l'engine generato in due giorni, ``senza guardarlo'', sia
``molto meglio di Godot'' nel 3D.

\paragraph{La previsione.} Il docente si aspetta quindi che ``a brevissimo''
succeda una di due cose:
\begin{itemize}
  \item esca un \textbf{engine pubblico gratuito} (open source) di questa nuova
        generazione, e ci si sposti su quello;
  \item oppure ogni studio si faccia \textbf{il suo}: ``mi serve un gioco 3D
        che ha tante funzionalità, che ne so, di musica'', ci si fa fare
        l'engine e poi ci si costruisce il gioco sopra. ``Probabile, molto
        probabile.''
\end{itemize}

\paragraph{Approfondimento: l'engine come scelta di design.} Jesse Schell, in
\emph{The Art of Game Design} (capitolo dedicato al prototipo, \emph{Prototyping
Tip \#8: Pick a ``Fast Loop'' Game Engine}), ricorda che la scelta dell'engine
non è solo tecnica: conta soprattutto \textbf{quanti cicli di prova al giorno}
permette. Lo sviluppo tradizionale è ``come cuocere il pane'' (scrivi il codice,
compili, avvii, raggiungi il punto da testare, provi, e se non va ricominci
tutto), mentre un engine con un buon sistema di scripting permette di
cambiare il codice mentre il gioco è in esecuzione, ``come lavorare
l'argilla''. Schell cita proprio Unreal, che sceglie ``un sistema di scripting
drag and drop (i Blueprint) che lavora insieme alla programmazione in C++''. Il
discorso del docente aggiunge un criterio nuovo a quello di Schell: oltre alla
velocità del ciclo, il peso di tutto ciò che l'engine porta con sé e che il
gioco non usa.

% ---------------------------------------------------------------------
\section{Il punto sui task e su GitHub}
\label{gd16:sec:task-github}

Il docente controlla se si sia collegato qualcun altro da remoto e poi
apre la lezione vera e propria con la domanda di rito: ``come eravate messi
con i task?''. Alcuni studenti gli hanno già scritto in privato (``quindi
bene''), ma vuole sapere ``in generale'' se ci sono difficoltà. Le risposte
sono rassicuranti (``ok, ok''): ``allora più avanti li vediamo''.

Una domanda riguarda \textbf{GitHub}, la lezione del
capitolo~\ref{gd13:cap:git-componenti}: qualcuno non c'era. La risposta è
tranquilla ma ferma: ``tranquilli. \textbf{Recuperatela}, perché è molto
importante'' (dal prossimo progetto, di gruppo, le consegne passeranno di
lì).

% ---------------------------------------------------------------------
\section{Il template in prima persona}
\label{gd16:sec:template}

\subsection{Perché proprio questo}

Per la lezione si crea un progetto nuovo dal \textbf{template in prima
persona} (\emph{First Person}). Non c'è un motivo legato ai contenuti, precisa il
docente (``non perché ci siano [sparatutto] nostri''): ``semplicemente perché
non l'abbiamo ancora visto''. Nel capitolo~\ref{gd12:cap:array-timeline} si
era sempre partiti dal template in terza persona, con la promessa di vedere più
avanti anche quello in prima persona e quello del veicolo.

Il template in prima persona ``è come quello in terza, ma in prima'': lo
stesso impianto, con qualche colore diverso.\footnote{Nella trascrizione
``Sesso. Cambio colore'': quasi certamente ``stesso, cambia colore''.} ``Poi c'è
il personaggio un po' \emph{clippato}'', cioè con parti della mesh che
attraversano la telecamera o altre geometrie, ``però vabbè, chi se ne
frega''. (Più avanti, durante l'esercitazione, il docente spiegherà
perché nel personaggio in prima persona si vedono due modelli compenetrati,
sezione~\ref{gd16:sec:suggerimenti}.)

\subsection{Se il template non c'è: installarlo dal launcher}

A qualcuno il template non compare. Non è un errore del progetto: i template
sono contenuti \textbf{opzionali} dell'installazione dell'engine, e si
aggiungono dal launcher.
\begin{enumerate}
  \item Nell'\textbf{Epic Games Launcher}, sezione di Unreal Engine, accanto
        alla versione installata si apre il menu della versione e si sceglie
        la voce per \textbf{gestire i moduli}\footnote{Nella trascrizione
        ``gestisci moduli''; nella versione inglese del launcher la voce
        corrispondente è \emph{Options}, nel menu a tendina accanto al pulsante
        di avvio della versione.};
  \item si spunta \textbf{Template e pacchetti di funzionalità}
        (\emph{Templates and Feature Packs});
  \item in fondo si preme \textbf{Applica}: il launcher scarica ``un giga'' e
        lo installa.
\end{enumerate}

\paragraph{Risparmiare spazio.} La stessa schermata, nota il docente, ``è super
comoda per \textbf{risparmiare spazio}'', perché permette di installare
\textbf{soltanto le cose che servono}. Alcuni componenti probabilmente ``non vi
serviranno mai, o a breve, comunque''; per esempio il supporto per la
\textbf{TV di iOS} (\emph{tvOS}): ``non mi è mai servito''. Altri, invece,
``potrebbero essere utili''. Vale la pena darci un'occhiata quando il disco
comincia a riempirsi: un'installazione completa di Unreal occupa molte decine
di gigabyte.


% ---------------------------------------------------------------------
\section{Non esiste un ``tasto salva tutto''}
\label{gd16:sec:perche}

``Quest'oggi volevo partire da una roba un pochettino concettuale, ma in realtà
torna molto utile per tipo tutto'': i \textbf{salvataggi}, cioè salvare
\textbf{proprio su disco} la partita.

Il docente parte smontando un'idea ingenua che molti hanno senza esserne
consapevoli:
\begin{quote}
Non so se ci avete mai pensato, ma in realtà non è che c'è un tasto salva, lui
salva tutto e poi fai ricarica e lui ricarica tutto. Non funziona così.
\end{quote}
In un engine non esiste una ``fotografia'' automatica del mondo di gioco da
mettere da parte e ripristinare. Il programmatore deve \textbf{dire
all'engine le singole variabili} da salvare, e \textbf{in che modo}. Tutto ciò
che non viene esplicitamente messo in un salvataggio, alla riapertura del gioco
semplicemente non c'è più.

\paragraph{Approfondimento: salvare è una decisione di design.} Il fatto che
il salvataggio vada costruito a mano ha una conseguenza che va oltre la
tecnica: \emph{che cosa} si salva, \emph{quando} e \emph{se} lo si può fare
sono decisioni di design. Jesse Schell (\emph{The Art of Game Design},
capitolo sul bilanciamento, a proposito delle punizioni) ricorda il caso di
\emph{Ultima~III} di Richard Garriott: un gioco da quasi cento ore in cui, se
morivano tutti e quattro i personaggi, il salvataggio veniva
\textbf{cancellato} e si ricominciava da capo. I giocatori lo trovavano
ingiusto, e presero l'abitudine di spegnere il computer quando i personaggi
erano in fin di vita, prima che il gioco facesse in tempo a cancellare il file:
``evitando di fatto la punizione''. Nello stesso libro (capitolo sulla storia)
Schell osserva che il salvataggio è una sorta di \emph{macchina del tempo}
messa in mano a ogni protagonista: tutto ciò che di grave accade può essere
annullato ricaricando, e questo rende difficilissimo raccontare una vera
tragedia in un videogioco. Chi progetta il sistema di salvataggio, quindi,
progetta anche quanto pesano le conseguenze delle scelte del giocatore.

% ---------------------------------------------------------------------
\section{La classe Save Game}
\label{gd16:sec:savegame}

\subsection{Creare il salvataggio}

Il primo passo è creare la \textbf{classe di salvataggio}:
\begin{enumerate}
  \item nel Content Browser, tasto destro, si crea una cartella
        \textbf{Save};
  \item dentro, tasto destro, \textbf{Blueprint Class}\footnote{Nella
        trascrizione ``Bluetooth class''.}, e si cerca come classe madre
        \textbf{SaveGame};
  \item il docente la chiama \textbf{PlayerSaves}.
\end{enumerate}

\paragraph{Quanti salvataggi?} Le classi di tipo Save Game ``sono proprio i
\textbf{tipi} di salvataggio''. Tendenzialmente in un gioco ce ne sono
\textbf{una o due}:
\begin{itemize}
  \item una per la \textbf{partita} (``salva la partita'');
  \item una per le \textbf{impostazioni} (volume, risoluzione, comandi: tutto
        ciò che vale a prescindere dalla partita in corso), ``perché anche le
        impostazioni vanno salvate, chiaramente''.
\end{itemize}
Qui ne basta una, per la partita.

\subsection{Un Blueprint senza viewport}

Aperto, \emph{PlayerSaves} ``è un Blueprint tipo quelli che abbiamo già visto, ma
un pochettino strano'': \textbf{non ha una parte visiva}, una viewport.
``Non gli serve'': essendo un salvataggio, non esiste nel mondo, quindi non ha
senso dargli una forma o una posizione. È parente, in questo, degli Actor
Component e della Game Instance del capitolo~\ref{gd13:cap:git-componenti}.

\paragraph{L'unico modo, a parte il C++.} Il docente sottolinea il ruolo di
questa classe: ``questo è l'\textbf{unico modo}, a parte C++, per salvare cose
su disco effettivamente in modo \textbf{permanente}''.

\paragraph{Contano solo le variabili.} Del Save Game ``non ci interessa questo
grafico, non ci interessano le funzioni, nulla'': interessano \textbf{soltanto
le variabili}, ``che sono quelle che vengono salvate''. Si crea quindi una
variabile:
\begin{itemize}
  \item nome: \textbf{Score}\footnote{Nella trascrizione ``Scorra''; poco dopo il
        docente parla di ``getScore'' e ``set dello score''.};
  \item tipo: \textbf{Integer}, ``mettiamo che è il punteggio del nostro gioco''.
\end{itemize}
Questa variabile, alla fine, verrà scritta sul disco. Ma il salvataggio, da
solo, non fa niente: va \textbf{creato}, \textbf{caricato} e
\textbf{salvato}, e ``ci sono un tot di step in mezzo''.

% ---------------------------------------------------------------------
\section{La Game Instance, rivista}
\label{gd16:sec:game-instance}

\subsection{Perché il salvataggio si gestisce lì}

Dove mettere il codice che crea e carica il salvataggio? Non su una mappa:
tutte queste cose ``vengono gestite al di fuori della mappa stessa''. Se le si
mettesse su un livello, ``ricaricherebbe il salvataggio per ogni mappa che
cambiamo'', e ``sarebbe parecchio complicato''. Il posto giusto è la
\textbf{Game Instance}, ``accennata la volta scorsa''
(sezione~\ref{gd13:sec:game-instance}) e che ``adesso chiaramente vediamo un
pochettino meglio''.

Si crea quindi (Blueprint Class, \emph{GameInstance}) una Game Instance, che il
docente chiama \textbf{MyGameInstance}. Il promemoria è lo stesso del capitolo
precedente:
\begin{itemize}
  \item ce n'è \textbf{solamente una} per ogni gioco;
  \item è \textbf{la prima cosa che viene caricata} e \textbf{l'ultima che viene
        scaricata}: ``è sempre lì, insomma''.
\end{itemize}

\subsection{Che cosa ci va, oltre ai salvataggi}

Uno studente chiede se la Game Instance serva solo per salvare. ``No'': ci
vanno tutte le cose ``gestite al di fuori'' dei livelli, quelle ``più
\textbf{globali}'':
\begin{itemize}
  \item i \textbf{salvataggi};
  \item gli \textbf{achievement}, ``queste cose astratte che non sono legate al
        livello stesso'';
  \item a volte perfino \textbf{la musica}, ``perché vuoi magari che non si senta
        il taglio tra un livello e l'altro'': una musica avviata da un Actor del
        livello verrebbe interrotta al cambio di mappa.
\end{itemize}

\subsection{Perché sopravvive ai cambi di mappa}

Il motivo è sempre lo stesso, e il docente lo ripete con un esempio
concreto:
\begin{quote}
Quando io faccio ``mi apri questa mappa'', lui distrugge tutto e ricrea tutto
sempre, ok? Tranne la Game Instance.
\end{quote}
Se nella Game Instance si mette una variabile intera per le monete e la si
porta da 0 a 25 raccogliendole, al cambio di mappa resta a \textbf{25}; ``tutto
il resto tornerebbe a zero''. (Una variabile così, precisa, qui ``non ci serve'':
era solo un esempio. Quello che interessa adesso è usare la Game Instance per un
salvataggio \emph{permanente}.)

La figura~\ref{gd16:fig:vita-dati} mette in fila i tre ``livelli di memoria''
visti finora: le variabili degli Actor, che muoiono a ogni cambio di livello;
quelle della Game Instance, che vivono finché il gioco è aperto; quelle del
Save Game scritte sul disco, che sopravvivono anche alla chiusura.

\begin{figure}[htbp]
\centering
\resizebox{\linewidth}{!}{%
\begin{tikzpicture}[
  gd16bar/.style={rounded corners=2pt, minimum height=0.62cm, inner sep=2pt,
                  font=\sffamily\scriptsize, anchor=west, text=black},
  gd16lab/.style={font=\sffamily\footnotesize, anchor=east, align=right},
  gd16ev/.style={font=\sffamily\scriptsize, align=center},
  x=1cm, y=1cm]
  % assi
  \draw[-{Stealth[length=2.5mm]}, thick] (0,0) -- (13.2,0)
       node[right, font=\sffamily\scriptsize] {tempo};
  % eventi
  \foreach \x/\t in {0.3/avvio\\del gioco, 3.6/Open\\Level, 6.9/Open\\Level,
                     9.3/chiusura\\del gioco, 11.2/nuovo\\avvio}{
    \draw[dashed, gray] (\x,0) -- (\x,4.1);
    \node[gd16ev, below] at (\x,-0.05) {\t};
  }
  % actor del livello
  \node[gd16lab] at (-0.2,3.5) {Variabili di un Actor\\(monete in un livello)};
  \fill[red!25] (0.3,3.2) rectangle (3.6,3.8);
  \fill[red!25] (3.6,3.2) rectangle (6.9,3.8);
  \fill[red!25] (6.9,3.2) rectangle (9.3,3.8);
  \fill[red!25] (11.2,3.2) rectangle (13,3.8);
  \foreach \x in {3.6,6.9}{\draw[red!70!black, line width=1.2pt] (\x,3.15) -- (\x,3.85);}
  \node[font=\sffamily\scriptsize] at (1.95,3.5) {0 \,$\to$\, 25};
  \node[font=\sffamily\scriptsize] at (5.25,3.5) {0 (si azzera)};
  \node[font=\sffamily\scriptsize] at (8.1,3.5) {0};
  \node[font=\sffamily\scriptsize] at (12.1,3.5) {0};
  % game instance
  \node[gd16lab] at (-0.2,2.3) {Variabili della\\Game Instance};
  \fill[blue!20] (0.3,2.0) rectangle (9.3,2.6);
  \fill[blue!20] (11.2,2.0) rectangle (13,2.6);
  \node[font=\sffamily\scriptsize] at (4.8,2.3) {25 \,$\to$\, 25 \,$\to$\, 25 (Event Init all'avvio)};
  \node[font=\sffamily\scriptsize] at (12.1,2.3) {0};
  % save game
  \node[gd16lab] at (-0.2,1.1) {Save Game\\scritto sul disco};
  \fill[green!25] (0.3,0.8) rectangle (13,1.4);
  \node[font=\sffamily\scriptsize] at (6.6,1.1) {il file \texttt{.sav} resta: al nuovo avvio lo si ricarica};
\end{tikzpicture}}
\caption{Quanto vive un dato. Un Open Level distrugge e ricrea tutti gli Actor
(e le loro variabili); la Game Instance sopravvive ai cambi di livello ma non
alla chiusura del gioco; solo ciò che è stato scritto con un Save Game resta
sul disco e può essere ricaricato al successivo avvio.}
\label{gd16:fig:vita-dati}
\end{figure}

\subsection{Event Init}

Nel grafo della Game Instance c'è l'\textbf{Event Init}, ``che sarebbe come
il Begin Play, uguale identico, ma per la Game Instance''. ``È la prima cosa
che viene chiamata quando parte il gioco, fondamentalmente.'' \emph{Init},
``che sta per \emph{initialize}, credo''. È qui che si scriverà il codice che
prepara il salvataggio.

% ---------------------------------------------------------------------
\section{I nodi del salvataggio}
\label{gd16:sec:nodi}

\subsection{Che cosa c'è}

Trascinando un filo dall'Event Init e cercando \textbf{save}, compaiono
``tutte quante le robe di salvataggio''. Il docente le zooma (``altrimenti, se
faccio tasto destro, non avrete zoomato il testo'') e le elenca:
\begin{itemize}
  \item un nodo per \textbf{cancellare} i salvataggi (\emph{Delete Game in Slot});
  \item uno per \textbf{crearli} (\emph{Create Save Game Object});
  \item uno per \textbf{caricarli} (\emph{Load Game from Slot});
  \item uno per \textbf{controllare se ce n'è già uno} (\emph{Does Save Game
        Exist});
  \item e quello per \textbf{salvarli} davvero (\emph{Save Game to Slot}), che
        si userà più avanti.
\end{itemize}
Prima di sceglierli, il docente chiede alla classe che cosa servirebbe
all'avvio del gioco, ``ammesso che non abbiamo la classica selezione del
salvataggio'' e che il gioco ``faccia già tutto lui''. Uno studente propone di
caricare l'ultimo salvataggio: ``sì, quello lì servirà, però dobbiamo vedere un
po' di cose prima''.

\subsection{Nodi istantanei e nodi asincroni}

Alcuni nodi hanno una \textbf{variante async}. ``Tutte quante quelle che trovate
con \emph{async} vuol dire che sono \textbf{asincrone}.'' Per capire che cosa
significhi, il docente ricorda come vengono eseguiti i nodi:
\begin{itemize}
  \item quasi tutti i nodi di Unreal\footnote{Nella trascrizione ``tutti quanti
        i nodi di Arrio'' (e poco più avanti ``i meccanismi di Ariadne''):
        storpiature di Unreal.}, ``tranne alcuni che avete già visto'', si
        eseguono \textbf{nello stesso frame} in cui vengono chiamati: ``partono,
        finiscono, tutto a posto. Tipo, fate 5 più 6, funziona subito'';
  \item altri invece \textbf{vanno nel tempo}: le \textbf{timeline}
        (sezione~\ref{gd12:sec:timeline}), i \textbf{Delay}. Questi hanno
        ``tendenzialmente'' un \textbf{orologino} disegnato sul nodo: ``se metto
        un delay, vedete che ha l'orologino qua''. (Sono i nodi
        \emph{latenti}, che per questo non si possono mettere dentro le
        funzioni, sezione~\ref{gd13:sec:funzioni}.)
\end{itemize}

\paragraph{A che cosa serve salvare in modo asincrono.} Si pensi a un gioco
grande, ``tipo un \textbf{open world survival}'', con tantissime cose da
salvare, magari con l'\textbf{autosave}. Se il salvataggio fosse istantaneo nel
senso di \emph{bloccante}, ``si bloccherebbe il gioco finché non ha finito, e
poi continuerebbe dopo''. Il salvataggio asincrono invece lavora \textbf{mentre
il gioco va}:
\begin{quote}
Sapete quando vi lampeggia l'icona del salvataggio in basso a destra? Quella
roba lì. Quindi ci mette del tempo per salvare.
\end{quote}
Qui però non serve: ``facciamo un gioco semplice, che salva quando facciamo
salva, quando finisce un livello''. La differenza è schematizzata nella
figura~\ref{gd16:fig:async}.

\begin{figure}[htbp]
\centering
\begin{tikzpicture}[x=0.62cm, y=1cm,
  gd16fr/.style={draw=black!60, fill=blue!12, minimum width=0.56cm,
                 minimum height=0.5cm, inner sep=0pt},
  gd16frb/.style={draw=black!60, fill=red!30, minimum height=0.5cm,
                  inner sep=0pt},
  gd16t/.style={font=\sffamily\footnotesize, anchor=east, align=right}]
  % sincrono
  \node[gd16t] at (-0.4,1.6) {Save Game to Slot\\(sincrono)};
  \foreach \i in {0,...,3}{\node[gd16fr] at (\i+0.5,1.6) {};}
  \node[gd16frb, minimum width=4.28cm, font=\sffamily\tiny] at (7.5,1.6)
       {un frame lunghissimo: gioco fermo};
  \foreach \i in {11,...,14}{\node[gd16fr] at (\i+0.5,1.6) {};}
  % asincrono
  \node[gd16t] at (-0.4,0.3) {Async Save Game\\to Slot};
  \foreach \i in {0,...,14}{\node[gd16fr] at (\i+0.5,0.3) {};}
  \draw[thick] (4,-0.05) -- (4,-0.2) -- (11,-0.2)
       node[midway, below=2pt, font=\sffamily\scriptsize]
       {il salvataggio prosegue in sottofondo (l'icona che lampeggia); poi scatta \emph{Completed}} -- (11,-0.05);
  \draw[-{Stealth}, thick] (0,2.35) -- (15.6,2.35)
       node[right, font=\sffamily\scriptsize] {frame};
\end{tikzpicture}
\caption{Salvataggio sincrono e asincrono. Il nodo normale svolge tutto il
lavoro dentro un solo frame, che con molti dati diventa un blocco visibile; la
variante \emph{async} lo distribuisce nel tempo, come una timeline o un Delay,
e segnala la fine con un pin d'uscita apposito.}
\label{gd16:fig:async}
\end{figure}

\paragraph{Approfondimento: i nodi asincroni nel dettaglio.} Le due varianti
asincrone si chiamano \emph{Async Save Game to Slot} e \emph{Async Load Game
from Slot}. Come gli altri nodi latenti hanno, oltre al pin d'esecuzione
normale (che prosegue subito), un pin \textbf{Completed} che scatta solo quando
l'operazione è finita, con un booleano \emph{Success} per il salvataggio o il
Save Game caricato per il caricamento. Il codice che ha bisogno del risultato
va attaccato a \emph{Completed}, non all'uscita normale: è l'errore tipico di
chi le usa la prima volta.

% ---------------------------------------------------------------------
\section{All'avvio: controlla, crea o carica}
\label{gd16:sec:init}

\subsection{Il primo avvio: creare}

``Tra crea, controlla, cancella, carica: cosa potrebbe servirci quando apriamo
il gioco per la prima volta, magari?'' La risposta della classe è
\textbf{creare}: c'è \textbf{Create Save Game Object}, che ``crea un salvataggio
nuovo''. Nel suo parametro \textbf{Save Game Class} bisogna dirgli quale
salvataggio, ``ed è il nostro, quello che abbiamo chiamato qui'': PlayerSaves.

\subsection{Ma forse c'è già: controllare}

Però non è detto che sia la prima volta: ``magari c'è già un salvataggio''. Per
questo c'è \textbf{Does Save Game Exist}, che controlla se esista già. Lo si
mette subito dopo l'Event Init.

\paragraph{Lo slot name.} Il parametro \textbf{Slot Name} è ``il nome proprio del
file che mettete sul disco''. Il docente lo chiama \textbf{Save1}\footnote{Nella
trascrizione ``chiamiamolo DOM. Save 1''; più avanti è chiaro che il valore usato
è ``save 1'', senza spazi o con, ``non è importante''.} (``non è importante'' quale
nome si scelga). Il nodo restituisce un \textbf{booleano}: vero se il salvataggio
esiste già, falso se non esiste. Da lì un \textbf{Branch}:
\begin{itemize}
  \item sul \textbf{False}: \textbf{Create Save Game Object};
  \item sul \textbf{True}: \textbf{Load Game from Slot}\footnote{Nella
        trascrizione ``lot save, load game from flot''.}, ``ecco qua,
        fantastico''.
\end{itemize}

\subsection{Un solo nome per tutti gli slot}

``Mi raccomando'': il nome dello slot deve essere \textbf{sempre lo stesso} in
tutti i nodi, ``perché altrimenti poi chiaramente non funziona''. Due strade:
\begin{itemize}
  \item copiarlo e incollarlo ovunque;
  \item oppure, ``per comodità'', creare una \textbf{variabile String} con il nome
        del salvataggio, \textbf{SaveName}, valore \emph{Save1}, e collegarla a
        tutti i pin \emph{Slot Name}. ``Così, se cambiamo questo qua di base,
        cambierà anche gli altri di conseguenza.''
\end{itemize}
E ancora: ``mi raccomando, nella variabile \textbf{scrivete qualcosa}, perché se
lo lasciate vuoto non può salvare un file senza nome''.

% ---------------------------------------------------------------------
\section{La reference al salvataggio}
\label{gd16:sec:reference}

\subsection{Tenersi il salvataggio in una variabile}

Sia \emph{Create Save Game Object} sia \emph{Load Game from Slot} restituiscono,
nel pin \textbf{Return Value}, un oggetto salvataggio: ``questo PlayerSaves che è
quello che vogliamo noi, fondamentalmente''. Va conservato, perché servirà dopo
``per metterlo sul counter, eccetera, per salvarlo un'altra volta, per un tot di
cose''. Si crea quindi una variabile per ``mantenere la \textbf{reference}, un
riferimento al nostro salvataggio'', chiamata \textbf{Save}.

\paragraph{Il tipo della variabile è una classe.} ``E qua si entra effettivamente
nei meccanismi di Unreal un attimino più completi'': il tipo della variabile non
è un numero né un testo, ma \textbf{il nostro salvataggio}: si cerca
\emph{PlayerSaves} tra i tipi (\emph{Object Reference}). ``Così, accedendo a
questo qua, troveremo tutte quante le informazioni che c'erano qua'', cioè tutte
le variabili definite in PlayerSaves.

Il docente sa che ``può sembrare un attimino complicata all'inizio'', e la
riconduce a una cosa già vista: quando si fa un \emph{Cast To
FirstPersonCharacter} dopo un overlap con una trigger box, il pin d'uscita
``è del tipo del personaggio''. Allo stesso modo una variabile può essere ``del
tipo del salvataggio''. In generale, ``qualsiasi cosa voi creiate qui dentro
praticamente può essere una variabile'': ogni Blueprint creato nel progetto
diventa anche un \textbf{tipo} di variabile, ``e così possiamo gestire tutte le
varie eventualità''.

\subsection{Il ramo del Create e il ramo del Load}

\begin{itemize}
  \item Sul ramo della \textbf{creazione} è semplice: il Return Value di
        \emph{Create Save Game Object} si collega direttamente a un
        \textbf{Set Save}. ``Creiamo il salvataggio e lo impostiamo così.''
  \item Sul ramo del \textbf{caricamento}, ``purtroppo'', non si può fare lo
        stesso, ``o almeno così facilmente''. \emph{Load Game from Slot}
        restituisce un Save Game \emph{generico}, perché ``potrebbero esserci
        diverse classi di salvataggio'': ``abbiamo creato questo qua, ma potrebbe
        essercene più di una''. Serve quindi un \textbf{Cast To PlayerSaves},
        e solo il suo pin d'uscita si collega al \textbf{Set Save}.
\end{itemize}
Il risultato, riassume il docente, è che ``qua in fondo siamo sicuri di avere
la variabile del salvataggio, che avrà al suo interno altre variabili,
chiaramente'': se il salvataggio non esiste lo crea e se lo tiene; se esiste lo
carica e se lo tiene (figura~\ref{gd16:fig:init}).

\begin{figure}[htbp]
\centering
\resizebox{\linewidth}{!}{%
\begin{bpgraph}
  \bpnode{ei}{0,0}{bpHEvent}{Event Init}{2.2}{}{exec/}
  \bpget{sn1}{0.2,-1.5}{str}{SaveName}{2.0}
  \bpnode{ex}{3.0,0}{bpHFunc}{Does Save Game Exist}{3.4}{exec/, str/Slot Name, int/User Index \bpval{0}}{exec/, bool/Return Value}
  \bpnode{br}{7.2,0}{bpHFlow}{Branch}{2.1}{exec/, bool/Condition}{exec/True, exec/False}
  \bpget{sn2}{8.0,-2.0}{str}{SaveName}{2.0}
  \bpnode{lo}{10.4,0.4}{bpHFunc}{Load Game from Slot}{3.4}{exec/, str/Slot Name, int/User Index \bpval{0}}{exec/, obj/Return Value}
  \bpnode{ca}{14.6,0.4}{bpHCast}{Cast To PlayerSaves}{3.6}{exec/, obj/Object}{exec/, exec/Cast Failed, obj/As Player Saves}
  \bpnode{s1}{19.0,0.4}{bpHVar}{SET Save}{2.4}{exec/, obj/Save}{exec/, obj/}
  \bpnode{cr}{10.4,-3.0}{bpHFunc}{Create Save Game Object}{5.4}{exec/, class/Save Game Class \bpval{PlayerSaves}}{exec/, obj/Return Value}
  \bpnode{s2}{19.0,-3.0}{bpHVar}{SET Save}{2.4}{exec/, obj/Save}{exec/, obj/}
  \bpwire{exec}{ei-o1}{ex-i1}
  \bpwire{str}{sn1-o1}{ex-i2}
  \bpwire{exec}{ex-o1}{br-i1}
  \bpwire{bool}{ex-o2}{br-i2}
  \bpwire{exec}{br-o1}{lo-i1}
  \bpwire{str}{sn2-o1}{lo-i2}
  \bpwire{exec}{br-o2}{cr-i1}
  \bpwire{exec}{lo-o1}{ca-i1}
  \bpwire{obj}{lo-o2}{ca-i2}
  \bpwire{exec}{ca-o1}{s1-i1}
  \bpwire{obj}{ca-o3}{s1-i2}
  \bpwire{exec}{cr-o1}{s2-i1}
  \bpwire{obj}{cr-o2}{s2-i2}
  \bpnote{14.6,-1.6}{il Load restituisce un Save Game generico:\\serve il cast al nostro tipo}
  \bpnote{3.0,-2.15}{lo stesso nome di slot,\\da una sola variabile String}
\end{bpgraph}}
\caption{La Game Instance all'avvio. \emph{Does Save Game Exist} controlla se
sul disco c'è già un file con quel nome di slot: se c'è, lo si carica e lo si
converte con un cast in PlayerSaves; se non c'è, se ne crea uno nuovo. In
entrambi i casi l'oggetto finisce nella variabile \emph{Save}, la reference che
tutto il resto del gioco userà.}
\label{gd16:fig:init}
\end{figure}

\paragraph{Approfondimento: lo stesso schema nel libro di Romero e Sewell.}
\emph{Blueprints Visual Scripting for Unreal Engine 5} (capitolo~11, \emph{Game
States and Applying the Finishing Touches}) costruisce un sistema quasi
identico per salvare il \textbf{round} raggiunto dal giocatore in uno sparatutto
in prima persona: una classe figlia di \emph{SaveGame} (\emph{BP\_SaveInfo}) con
una sola variabile intera \emph{Round}; una variabile \emph{SaveInfoRef} di tipo
\emph{Object Reference} a quella classe; una variabile String
\emph{SaveSlotName} con il nome del file. Il caricamento (\emph{LoadRound}) è
lo stesso della figura~\ref{gd16:fig:init}: \emph{Does Save Game Exist},
\emph{Branch}, \emph{Load Game from Slot}, \emph{Cast To BP\_SaveInfo}, set
della reference. Due differenze utili:
\begin{itemize}
  \item il libro mette tutto nel \emph{FirstPersonCharacter}, dentro due
        \textbf{macro}; a lezione si è scelta la Game Instance, che ha il
        vantaggio di non dipendere dal personaggio e di non ricaricare nulla a
        ogni cambio di livello;
  \item nel salvataggio (\emph{SaveRound}) il libro usa un \textbf{Is Valid}
        sulla reference: se non è valida crea il Save Game, poi aggiorna la
        variabile e chiama \emph{Save Game to Slot}. È un controllo prudente
        da aggiungere se si teme che la reference sia vuota.
\end{itemize}
Il quiz di fine capitolo del libro fissa due punti visti anche a lezione: la
classe usata per i dati da salvare è \emph{SaveGame}, e \emph{Load Game from
Slot} restituisce un'istanza di SaveGame che va convertita con un \emph{Cast
To} per accedere alle variabili salvate. Il libro mostra anche il nodo per
\textbf{cancellare} il salvataggio, \emph{Delete Game in Slot}, usato per un
pulsante \emph{Reset All}: prima si controlla con \emph{Does Save Game Exist},
poi si cancella e si ricarica il livello.

% ---------------------------------------------------------------------
\section{Dirlo all'engine e salvare davvero}
\label{gd16:sec:salvare}

\subsection{Maps \& Modes}

``Già dovrebbe funzionare. L'unica cosa è che non gli abbiamo messo la Game
Instance ancora'', cioè non si è detto all'engine di usarla. Come nel capitolo
precedente: \textbf{Edit $\rightarrow$ Project Settings $\rightarrow$ Maps \&
Modes}, ``dove si impostano insomma tutte queste cose'', e alla voce
\textbf{Game Instance Class} si sceglie la propria, ``perché le altre sono
dell'engine di base, non ci interessano''. Uno studente non la trova nell'elenco:
``quella che hai creato tu. Non so come l'hai chiamata. Io l'ho chiamata
\emph{My}, per distinguere''.

\subsection{Creare non è salvare}

Ora, premendo Play, la Game Instance controlla se c'è il salvataggio. ``Per ora
fallirà sicuramente'', quindi ne creerà uno. Ma attenzione:
\begin{quote}
Quando fa create, non vuol dire che lo salva su disco: vuol dire soltanto che
lui se l'è creato temporaneamente.
\end{quote}
Il file sul disco nasce solo con l'operazione di \textbf{salva}, che va
chiamata ``quando effettivamente vogliamo salvare''.

\subsection{Save Game to Slot}

``Per fare ciò è abbastanza facile'': il nodo \textbf{Save Game to Slot} (c'è
anche la variante asincrona: ``scegliete voi quello che preferite''; la
variante, ricorda, ``ci può mettere più di un frame, perché ha la rotellina''). I
suoi parametri:
\begin{itemize}
  \item \textbf{Save Game Object}: l'oggetto da salvare, ``che ce l'abbiamo,
        perché lo impostiamo sempre qui'': il \emph{Get Save};
  \item \textbf{Slot Name}: ``serve sempre lo stesso'', quindi la variabile
        \emph{SaveName}.
\end{itemize}
``E così siamo a posto.'' A questo punto ci sono ``tutte quante le funzionalità
base di un salvataggio''.

\subsection{Il custom event SaveGame}

Manca un modo per \textbf{chiamare} il salvataggio, ``perché altrimenti non viene
mai chiamato'': il nodo non ha un evento d'inizio. Ce lo si crea: \textbf{Add
Custom Event}, chiamato \textbf{SaveGame}\footnote{Nella trascrizione ``Boxer
Game''; subito dopo il docente dice ``chiamiamo questo save game qui''.}.
\begin{quote}
Da qualsiasi punto del gioco chiamiamo questo SaveGame qui: vuol dire ``salva il
gioco, fai tu''. Così la prossima volta potrà caricarlo.
\end{quote}
E poiché vive nella Game Instance, è raggiungibile da qualunque Blueprint
(figura~\ref{gd16:fig:save}).

\begin{figure}[htbp]
\centering
\resizebox{0.78\linewidth}{!}{%
\begin{bpgraph}
  \bpnode{ce}{0,0}{bpHEvent}{Custom Event: SaveGame}{3.6}{}{exec/}
  \bpget{sv}{1.0,-1.5}{obj}{Save}{1.6}
  \bpget{sn}{1.0,-2.3}{str}{SaveName}{2.0}
  \bpnode{sg}{4.6,0}{bpHFunc}{Save Game to Slot}{3.6}{exec/, obj/Save Game Object, str/Slot Name, int/User Index \bpval{0}}{exec/, bool/Return Value}
  \bpwire{exec}{ce-o1}{sg-i1}
  \bpwire{obj}{sv-o1}{sg-i2}
  \bpwire{str}{sn-o1}{sg-i3}
\end{bpgraph}}
\caption{Il salvataggio vero e proprio, nella Game Instance: un custom event
\emph{SaveGame}, chiamabile da qualunque punto del gioco, scrive sul disco
l'oggetto tenuto nella variabile \emph{Save}, sempre con lo stesso nome di
slot.}
\label{gd16:fig:save}
\end{figure}

``Non penso che ci servano altre cose.'' Il docente aspetta che tutti abbiano
copiato (``confermatemi le cose, dovremmo avercelo'').


% ---------------------------------------------------------------------
\section{Le monete}
\label{gd16:sec:monete}

\subsection{Il Blueprint Coins}

Per avere qualcosa da salvare servono delle \textbf{monetine}. Nella cartella
del player: tasto destro, \textbf{Blueprint Class}, \textbf{Actor}, chiamato
\textbf{Coins}. Mentre lo si crea, il docente aiuta uno studente che non trova
la variabile del salvataggio nella Game Instance: ``si chiamerà \emph{Save} e
basta, perché prende il nome che gli avete messo''; se i nomi non coincidono,
``al massimo rinomina questo qua, così lo trovi più facilmente''.

\paragraph{La forma.} ``Su questa coin, cosa dobbiamo fare per raccoglierla?''
La classe propone una \textbf{box} di collisione; il docente suggerisce
``magari una sfera, che è più carina'', ma alla fine per la moneta si usa
semplicemente un \textbf{cilindro} (``eccola qui, fantastica'').

\paragraph{L'overlap e il cast.} Dall'evento di \textbf{overlap} della moneta
si fa, ``giusto per controllo'', un \textbf{Cast To FirstPersonCharacter}: ora
il personaggio non è più quello in terza persona, ``perché abbiamo questo qua
adesso''. L'esecuzione arriva lì ogni volta che il giocatore collide con una
moneta.

\subsection{Chi deve fare il lavoro: il personaggio}

Qui il docente introduce un principio di organizzazione del codice:
\begin{quote}
L'ideale sarebbe se tutto quanto il codice lo facesse direttamente il First
Person Character. Quindi non è che ogni monetina si salva la sua informazione,
ma lo fa il player, che diventa un pochettino più intuitivo.
\end{quote}
La moneta deve solo \emph{avvisare} di essere stata presa; che cosa succede al
punteggio e al salvataggio è affare del giocatore.

\paragraph{Il custom event CoinGrabbed.} Nel Blueprint del personaggio si
aggiunge quindi un \textbf{Add Custom Event} chiamato \textbf{CoinGrabbed}.
``Come vedete, questi custom event, una volta che vogliamo cominciare a
strutturare bene il tutto, si useranno molto spesso.'' Se ne può fare a meno
``la maggior parte delle volte'', ``tuttavia diventa un casino il codice,
semplicemente'': con i custom event ``potete riorganizzare un attimino meglio il
tutto''.

\paragraph{La moneta chiama il personaggio.} Sulla moneta, dal pin d'uscita del
cast (\emph{As First Person Character}) si trascina un filo e si cerca
\textbf{CoinGrabbed}: ``facilissimo''. Poi, ``sulla Coin non dobbiamo più fare
niente'', se non un \textbf{Destroy Actor}, ``così si distrugge la monetina
quando ci passiamo sopra'' (figura~\ref{gd16:fig:coin}).

A uno studente che non trova l'evento, il docente spiega da dove si prende:
lo si tira fuori dal pin del personaggio ottenuto con il cast, e appare solo
se l'evento è stato creato nel personaggio con quel nome. Il collegamento,
dice, fa ``proprio un \textbf{teletrasporto}'': quando l'esecuzione entra nel
nodo \emph{CoinGrabbed} sulla moneta, ``esce di qua'', cioè dal custom event
nel Blueprint del personaggio.

\begin{figure}[htbp]
\centering
\resizebox{\linewidth}{!}{%
\begin{bpgraph}
  \bpnode{ov}{0,0}{bpHEvent}{On Component Begin Overlap (Box)}{4.4}{}{exec/, obj/Other Actor}
  \bpnode{ca}{5.2,0}{bpHCast}{Cast To FirstPersonCharacter}{4.3}{exec/, obj/Object}{exec/, exec/Cast Failed, obj/As First Person Character}
  \bpnode{cg}{10.4,0}{bpHFunc}{CoinGrabbed}{2.6}{exec/, obj/Target}{exec/}
  \bpnode{da}{13.8,0}{bpHFunc}{Destroy Actor}{2.6}{exec/, obj/Target \bpval{self}}{exec/}
  \bpwire{exec}{ov-o1}{ca-i1}
  \bpwire{obj}{ov-o2}{ca-i2}
  \bpwire{exec}{ca-o1}{cg-i1}
  \bpwire{obj}{ca-o3}{cg-i2}
  \bpwire{exec}{cg-o1}{da-i1}
  \bpnote{10.4,-1.6}{l'evento vive nel personaggio:\\l'esecuzione ``si teletrasporta'' lì}
\end{bpgraph}}
\caption{Il Blueprint \emph{Coins}. La moneta controlla che a toccarla sia il
personaggio, gli comunica di essere stata presa chiamando il suo custom event
\emph{CoinGrabbed} e si distrugge. Tutta la logica del punteggio sta nel
personaggio, non nella moneta.}
\label{gd16:fig:coin}
\end{figure}

\paragraph{Prima prova.} Si piazzano un po' di monete ``qui in giro'', ci si
passa sopra e ``come vedete si distruggono'': il Destroy
Actor\footnote{Nella trascrizione ``destroy hardware''.} viene chiamato, e
quindi, ``immagino'', anche l'evento sul personaggio.

\paragraph{Attenzione all'ordine.} A uno studente la moneta viene raccolta ma
resta lì: il Destroy Actor va collegato dopo la chiamata, ``se no ti rimane lì
e quindi non la prendi mai, insomma. Cioè, la prendi, ma non si distrugge''.

% ---------------------------------------------------------------------
\section{La torre di nodi}
\label{gd16:sec:torre}

\subsection{Variabili dentro variabili}

``Ora da qua dobbiamo accedere a tutte quante le informazioni del salvataggio.''
Ed è ``la prima volta che vedremo per davvero come vengono fuori le cose con
Unreal: variabili con dentro variabili con dentro variabili''. Il percorso, nel
custom event \emph{CoinGrabbed} del personaggio, è questo:
\begin{enumerate}
  \item \textbf{Get Game Instance}: ``prendiamo la Game Instance semplicemente
        perché è lì che abbiamo messo i salvataggi'', e ce n'è solo una;
  \item \textbf{Cast To MyGameInstance}: ``dobbiamo controllare che sia il nostro
        tipo di Game Instance'' (lo stesso motivo del capitolo precedente: il
        nodo restituisce una Game Instance generica);
  \item \textbf{Get Save}: ``all'interno della Game Instance, se vi ricordate,
        abbiamo messo il salvataggio'';
  \item \textbf{Get Score}: ``all'interno del salvataggio c'è lo score delle
        monete''.
\end{enumerate}
``E quindi si creano sempre queste \textbf{referenze di referenze di
referenze}'', ``la nostra \textbf{torre di nodi}''.

\paragraph{Il cast puro.} ``Se vi dà fastidio questa roba qua'', cioè i pin
d'esecuzione del cast che allungano la catena, si può fare tasto destro sul
cast e \textbf{convertirlo in puro} (\emph{Convert to Pure Cast}), ``e funziona
comunque''. Un cast puro non ha pin d'esecuzione né il ramo \emph{Cast Failed}:
se fallisce restituisce semplicemente una reference vuota.

\subsection{Sommare uno, e il set}

``Tutto questo casino'' serve a ottenere il valore del punteggio, ``per
incrementarlo, chiaramente, quando prendiamo una moneta''. Il docente ricorda
come si fanno le somme, ``solo tanto per ricordarvelo'': si prende il valore,
si somma 1 e poi si fa il \textbf{Set} (come nel contatore delle morti,
sezione~\ref{gd13:sec:game-instance}). La differenza è solo che il set va fatto
\emph{sul salvataggio}: dal pin di \emph{Save} si prende non il \emph{Get Score}
ma il \textbf{Set Score}, ``che è la stessa roba, semplicemente collegato al
salvataggio''.

\paragraph{Funziona, ma non salva.} Con un \textbf{Print String} di prova, salendo
sulle monete compare 1, 2, 3. ``Fantastico. Però se riavvio, fa sempre 1, 2,
3'': il valore cambia nell'oggetto in memoria, ma non viene scritto da nessuna
parte. ``Perché non sto salvando, chiaramente.''

\subsection{Chiamare SaveGame}

L'ultimo pezzo è chiamare il custom event \textbf{SaveGame} della Game Instance,
``in modo che chiami questo e che salvi''. Lo si prende, ancora una volta, dal pin
della Game Instance convertita con il cast. ``So che è la prima volta che
interconnettiamo così tanto i Blueprint, ma poi vi verrà abbastanza naturale.''
Nel caso della lezione si salva \textbf{a ogni moneta presa}
(figura~\ref{gd16:fig:torre}).

\begin{figure}[htbp]
\centering
\resizebox{\linewidth}{!}{%
\begin{bpgraph}
  \bpnode{ce}{0,0}{bpHEvent}{Custom Event: CoinGrabbed}{3.8}{}{exec/}
  \bpnode{gi}{0,-2.2}{bpHPure}{Get Game Instance}{3.2}{}{obj/Return Value}
  \bpnode{ca}{4.6,0}{bpHCast}{Cast To MyGameInstance}{4.0}{exec/, obj/Object}{exec/, exec/Cast Failed, obj/As My Game Instance}
  \bpnode{gs}{9.4,-2.4}{bpHPure}{Get Save}{2.2}{obj/Target}{obj/}
  \bpnode{gsc}{12.0,-3.3}{bpHPure}{Get Score}{2.2}{obj/Target}{int/}
  \bpnode{add}{14.8,-2.9}{compact}{$+$}{1.3}{int/, int/\bpval{1}}{int/}
  \bpnode{ss}{16.8,0}{bpHFunc}{SET Score}{2.8}{exec/, obj/Target, int/Score}{exec/, int/}
  \bpnode{sg}{20.6,0}{bpHFunc}{SaveGame}{2.6}{exec/, obj/Target}{exec/}
  \bpwire{exec}{ce-o1}{ca-i1}
  \bpwire{obj}{gi-o1}{ca-i2}
  \bpwire{exec}{ca-o1}{ss-i1}
  \bpwire{obj}{ca-o3}{gs-i1}
  \bpwire{obj}{gs-o1}{gsc-i1}
  \bpwire{obj}{gs-o1}{ss-i2}
  \bpwire{int}{gsc-o1}{add-i1}
  \bpwire{int}{add-o1}{ss-i3}
  \bpwire{exec}{ss-o1}{sg-i1}
  \bpwire{obj}{ca-o3}{sg-i2}
  \bpnote{4.6,-3.4}{Game Instance $\rightarrow$ Save $\rightarrow$ Score:\\referenze di referenze}
  \bpnote{16.8,-1.9}{a lezione, tra SET e SaveGame,\\anche un Print String di prova\\del nuovo valore}
\end{bpgraph}}
\caption{La ``torre di nodi'' nel personaggio (senza il Print String di prova). A ogni moneta: si prende la Game
Instance, la si converte nel proprio tipo, si raggiunge il salvataggio e al suo
interno il punteggio, lo si incrementa di uno con il set, e infine si chiede
alla Game Instance di scrivere tutto sul disco con \emph{SaveGame}.}
\label{gd16:fig:torre}
\end{figure}

% ---------------------------------------------------------------------
\section{La prova, il file e il riepilogo}
\label{gd16:sec:prova}

\subsection{1, 2, 3\dots{} e poi 4, 5, 6}

Con il salvataggio collegato, ``funziona tutto'': si prendono tre monete, 1, 2,
3; si chiude e si riavvia, e si continua con 4, 5, 6; di nuovo, 7, 8, 9.
\begin{quote}
E se chiudessi il PC, lo spegnessi e lo riprendessi dopo, funzionerebbe sempre
così. Non è legato alla cosa.
\end{quote}
Il valore non dipende più dalla sessione di gioco né dall'editor aperto: è sul
disco.

\subsection{Dove finisce il file}

Il docente lo fa vedere ``nel pratico'', aprendo la cartella del progetto: in
\textbf{Saved $\rightarrow$ SaveGames} c'è il file \texttt{Save1.sav}, ``proprio
il file che è stato salvato su disco''. Lo slot name è, letteralmente, il nome
del file.
\begin{itemize}
  \item Nell'editor lo mette \textbf{nella cartella del progetto}, ``perché è
        come test'';
  \item quando si fa la \textbf{build}, lo mette ``nella cache del PC'', cioè
        in una cartella dell'utente.
\end{itemize}

\paragraph{Approfondimento: dove sono i salvataggi di un gioco pacchettizzato.}
Su Windows, un gioco Unreal pacchettizzato scrive di norma i salvataggi in
\texttt{\%LOCALAPPDATA\%\char92 <NomeProgetto>\char92
Saved\char92 SaveGames}, con un file \texttt{.sav} per ogni slot. È
utile saperlo per i playtest: cancellare quel file equivale a far ripartire il
gioco ``da nuovo'', ed è la stessa cosa che fa, dall'interno del gioco, il nodo
\emph{Delete Game in Slot}. Nell'editor basta cancellare il file in
\texttt{Saved/SaveGames} del progetto.

\subsection{Il riepilogo}

Il docente lascia finire di copiare (``se c'è qualcosa che non vi torna
concettualmente, ditemelo'') e poi riassume tutto ``molto rapidamente''
(figura~\ref{gd16:fig:mappa}):
\begin{enumerate}
  \item si è creata una classe \textbf{PlayerSaves}, il salvataggio: dentro ci
        sono tutte le variabili che vengono salvate, ``che è quello che viene
        salvato su disco'';
  \item quando il gioco inizia, la Game Instance \textbf{controlla} se c'è il
        salvataggio: se non c'è \textbf{lo crea}, se c'è \textbf{lo carica};
  \item quando si va sul personaggio, la moneta chiama \emph{CoinGrabbed}: si
        prende la Game Instance, si controlla che sia quella giusta, si prende il
        salvataggio, si prende lo score, ``perché tanto è qui dentro, appena
        caricato'', si somma 1 e si dice alla Game Instance di
        \textbf{salvare}. ``E basta, finito lì.''
\end{enumerate}

\begin{figure}[htbp]
\centering
\begin{tikzpicture}[
  gd16bp/.style={draw=black!70, rounded corners=3pt, fill=#1, align=center,
                 font=\sffamily\footnotesize, minimum width=3.4cm,
                 minimum height=1.25cm, inner sep=4pt},
  gd16ar/.style={-{Stealth[length=2.6mm]}, thick},
  gd16lb/.style={font=\sffamily\scriptsize, align=center, fill=white,
                 inner sep=1.5pt},
  node distance=1.2cm and 2.0cm]
  \node[gd16bp=yellow!20] (coin) {\textbf{Coins}\\(Actor, cilindro)};
  \node[gd16bp=orange!15, right=2.6cm of coin] (pl)
       {\textbf{FirstPersonCharacter}\\custom event CoinGrabbed};
  \node[gd16bp=blue!12, below=1.6cm of pl] (gi)
       {\textbf{MyGameInstance}\\Event Init, SaveGame,\\variabili Save e SaveName};
  \node[gd16bp=green!15, left=2.6cm of gi] (sv)
       {\textbf{PlayerSaves}\\(SaveGame)\\variabile Score};
  \node[draw=black!70, cylinder, shape border rotate=90, aspect=0.25,
        fill=gray!15, minimum width=2.2cm, minimum height=1.4cm,
        font=\sffamily\footnotesize, align=center, below=1.3cm of sv] (disk)
       {disco\\\texttt{Save1.sav}};
  \draw[gd16ar] (coin) -- node[gd16lb, above] {overlap $\rightarrow$ cast\\$\rightarrow$ CoinGrabbed} (pl);
  \draw[gd16ar] (pl) -- node[gd16lb, right] {Get Game Instance\\$\rightarrow$ cast $\rightarrow$ Get Save\\$\rightarrow$ Score $+1$ $\rightarrow$ SaveGame} (gi);
  \draw[gd16ar] (gi) -- node[gd16lb, above] {reference\\\emph{Save}} (sv);
  \draw[gd16ar] ([xshift=-6pt]sv.south) -- node[gd16lb, left] {Save Game\\to Slot} ([xshift=-6pt]disk.north);
  \draw[gd16ar] ([xshift=6pt]disk.north) -- node[gd16lb, right] {Load Game from Slot\\(all'avvio)} ([xshift=6pt]sv.south);
\end{tikzpicture}
\caption{Chi parla con chi. La moneta avvisa il personaggio; il personaggio
passa dalla Game Instance per raggiungere il salvataggio e modificarlo; la Game
Instance lo scrive sul disco (e all'avvio lo rilegge).}
\label{gd16:fig:mappa}
\end{figure}

\subsection{Salvare cose complesse}

La conseguenza generale: ``potete mettere qui dentro qualsiasi variabile e vi
salverà il gioco''. Ma, di nuovo, niente è automatico. Se si vuole salvare
qualcosa di complesso, ``tipo la \textbf{posizione del personaggio}'', bisogna
creare nel salvataggio un \textbf{Vector}, o ``addirittura una
\textbf{Transform}'' (posizione, rotazione e scala), e scrivere il codice che la
salva e che, al caricamento, \emph{rimette} il personaggio lì: ``però dovete
farlo voi, il codice è un attimino lungo''.

\paragraph{Approfondimento: decidere che cosa salvare.} Romero e Sewell, nel
progettare il loro salvataggio a round, fanno una scelta esplicita di design:
salvano \emph{solo} il round raggiunto e non quanti nemici sono stati uccisi,
perché ``avrebbe più senso per il giocatore'' che ogni sessione ricominci
dall'inizio di un round. È il punto da cui partire anche per il proprio
gioco: elencare che cosa il giocatore si aspetta di ritrovare (progressi,
inventario, posizione, impostazioni) e che cosa invece può ripartire da capo.
Più variabili si salvano, più codice serve per rimettere ogni cosa al suo posto
al caricamento.

% ---------------------------------------------------------------------
\section{Le domande della classe}
\label{gd16:sec:domande}

\paragraph{Da dove si prende SaveGame?} Uno studente non trova il nodo da
chiamare nel personaggio. ``Questo qua lo trascini fuori da qua, perché questo
SaveGame l'abbiamo creato noi'' nella Game Instance: ``fa parte della Game
Instance, quindi lo tiriamo fuori dalla Game Instance''. Prova ne è che, facendo
doppio clic sul nodo, si viene portati proprio lì; e la voce \textbf{Target} del
nodo ``ti indica dov'è questa cosa qui''.

\paragraph{Il numero non si incrementa.} A un altro il punteggio non sale tra
un avvio e l'altro: ``forse non lo carichi correttamente qua''. Il docente
guarda il collegamento: lo studente ha collegato ``il numero di coin'', ``fatti
vedere'', ma il numero di coin e lo score ``è la stessa roba''. Il problema era
altrove; dopo un controllo (``vediamo un attimo'', a un altro studente che ``non
sta salvando''), il sistema funziona: ``tutte quante le volte dovrebbe
incrementarti la monetina, infatti vedi? 12, lo faccio ripartire\dots{} 14, 15''.

\paragraph{Vedere il valore sempre a schermo.} Uno studente vorrebbe vedere il
punteggio a schermo in ogni momento, non solo quando prende una moneta. Senza
ancora i widget, ``purtroppo non puoi, tra virgolette'', ma c'è ``un piccolo
escamotage'': mettere il \textbf{Print String} sull'\textbf{Event Tick}, con
durata \textbf{zero}. A ogni frame il messaggio viene riscritto e sparisce
subito, ``e così si vede sempre, vedi? Così funziona''. (Il capitolo
precedente aveva mostrato un'alternativa, il plugin Text~3D.)

\paragraph{Quali nodi si possono rendere puri?} Dopo il cast puro, qualcuno
chiede se si possano convertire in puri anche gli altri nodi. ``No, no, no.
Soltanto i cast si può fare'' tra questi; ``gli altri, tendenzialmente, no''.
``Poi vedremo le cose pure, dopo.''

\paragraph{Il best score del prossimo task.} Prima di passare alla pratica il
docente chiede conferma che tutto funzioni a tutti, e ripete l'avviso: ``questa
cosa qui che viene salvato il \textbf{best score}, ve lo anticipo già, ci sarà
nel prossimo task. Quindi vi ricordate questa roba qua, abbastanza rapida anche
da ricopiare''.

\paragraph{Approfondimento: dal punteggio al best score.} Il passo dal
punteggio cumulativo della lezione al \emph{miglior} punteggio è piccolo ma
istruttivo: il punteggio della partita corrente non va salvato a ogni moneta
(deve ripartire da zero a ogni partita), mentre nel Save Game si tiene solo un
\emph{BestScore}. A fine partita si confronta: se il punteggio corrente è
maggiore del \emph{BestScore} salvato, lo si sostituisce e si chiama
\emph{SaveGame}; altrimenti non si scrive nulla. È lo stesso schema
\emph{get, confronta, set, salva} della figura~\ref{gd16:fig:torre}, con un
\textbf{Branch} su un confronto ($>$) al posto della somma.

% ---------------------------------------------------------------------
\section{L'esercitazione: un team, nove feature}
\label{gd16:sec:esercitazione}

\subsection{La simulazione}

``Facciamo un pochettino di pratica.''\footnote{Nella trascrizione ``Furnace un
pochettino di pratica''.} Il docente chiude il progetto dei salvataggi (``questo
qua non serve più'') e propone una simulazione:
\begin{quote}
Immaginate di star lavorando per un'azienda X, ok? Siete nello stesso team e
dovete implementare tot feature di test.
\end{quote}
Le regole:
\begin{itemize}
  \item le feature \textbf{si dividono tra i membri del team}: ``ne fate una a
        testa'', scegliendo ``quella che vi piace di più'', l'importante è che
        \textbf{non si faccia la stessa in più persone};
  \item chi finisce prima ``può fare le altre che avanzano'' (``appena ne finite
        una, potete passare a quella dopo, se era particolarmente facile'');
  \item si lavora \textbf{singolarmente}, ciascuno sul proprio progetto, ``per
        non allungare troppo'';
  \item vanno bene sia il template in prima sia quello in terza persona:
        ``funzionano tutti quanti sia con il first sia con il third''. Forse
        quella della torcia ``è un pochettino più bellina con il first, però
        funziona tranquillamente'';
  \item ognuno \textbf{scrive nella chat} quale feature ha scelto, ``così ce
        l'abbiamo tutti'' (dopo un momento di incertezza su quale chat usare:
        ``potete scrivere su questa chat o no? \dots{} Me lo metto
        sull'altro'');
  \item non è un task da consegnare: ``sì, sì, qua adesso, non è come task da
        fare in futuro''. E soprattutto: ``tranquilli, che \textbf{non è
        valutata} questa cosa qua. È soltanto per prenderci un attimino la mano
        con l'engine''. ``Vi do una mano tranquillamente. Non è che non mi dovete
        chiedere nulla.''
\end{itemize}

L'idea è la stessa del richiamo sulla partecipazione del
capitolo~\ref{gd13:cap:git-componenti}: lavorare come in un team, dividersi il
lavoro, chiedere.

\subsection{Le nove feature}

L'elenco dettato dal docente è riportato nella tabella~\ref{gd16:tab:feature}.

\begin{table}[htbp]
\centering
\small
\begin{tabular}{@{}r p{0.82\linewidth}@{}}
\toprule
\# & \textbf{Feature} \\
\midrule
1 & Un \textbf{Actor che si ingrandisce} quanto più è lontano dal player. \\
2 & Un \textbf{jetpack} per il player: si preme il tasto e si vola più in alto. \\
3 & Una \textbf{torcia}: si preme il tasto e si illumina; ``e poi magari si
    scarica anche una \textbf{batteria}''. \\
4 & Se il player \textbf{cade in una zona}, viene riportato a dove era prima che
    cadesse. \\
5 & Un oggetto messo nella scena che, tramite codice, \textbf{si vede soltanto il
    35\% delle volte} che si fa Play; le altre volte ``è come se non ci fosse mai
    stato''. \\
6 & Una \textbf{torretta} che spara proiettili che, quando colpiscono il
    giocatore, \textbf{riavviano il livello}. \\
7 & Un \textbf{tasto che genera un oggetto tre volte}; poi c'è un'attesa di
    \textbf{5 secondi}, dopo la quale se ne possono generare altri tre. \\
8 & Una \textbf{zona che toglie 10 HP al secondo}: ``quindi fate voi il sistema
    di HP''. \\
9 & Un \textbf{Actor che cambia modello a caso} ogni secondo. \\
\bottomrule
\end{tabular}
\caption{Le feature dell'esercitazione, una a testa. Il docente chiede di
sceglierne una, scriverla in chat e poi, finita, passare a una di quelle rimaste.}
\label{gd16:tab:feature}
\end{table}

\paragraph{Il 35\% e la generazione procedurale.} Sulla feature~5 il docente si
sofferma: ``serve molto, questo qua, per la \textbf{generazione procedurale}
delle varie cose''. Il requisito vero è sottile: ``dovete fare in modo che il
player \textbf{non si accorga} che c'era questa cosa'', almeno nelle volte in cui
non deve esserci. Non basta far sparire l'oggetto: deve sembrare che non sia mai
stato lì.

\subsection{Possibili strade}

La registrazione del lavoro individuale non è utilizzabile, quindi non si sa
come ciascuno abbia risolto la propria feature. Il riquadro che segue propone,
per ognuna, una strada con i nodi visti finora (o con pochi nodi nuovi, indicati
per nome): \textbf{è un'integrazione}, non quanto detto a lezione.

\paragraph{Approfondimento: una possibile soluzione per ciascuna feature.}
\begin{enumerate}
  \item \textbf{Scala con la distanza.} Sull'\emph{Event Tick} dell'Actor:
        \emph{Get Player Pawn} (sezione~\ref{gd16:sec:suggerimenti}),
        \emph{Get Distance To}, moltiplicato per un fattore (per esempio
        $0{,}002$) e magari limitato con un \emph{Clamp}; il risultato in
        \emph{Set Actor Scale 3D} (un float collegato a un vettore si converte
        da solo in $(s,s,s)$).
  \item \textbf{Jetpack.} Un input nel personaggio: finché il tasto è premuto,
        \emph{Launch Character} con una piccola velocità verso l'alto (o
        \emph{Add Movement Input} lungo $z$ in modalità di volo). Con un float
        \emph{Fuel} che scende con i \emph{Delta Seconds} si ottiene anche il
        limite di carburante.
  \item \textbf{Torcia con batteria.} Uno \emph{Spot Light} come componente figlio
        della telecamera; il tasto fa \emph{Toggle Visibility}. Un float
        \emph{Battery} cala sull'\emph{Event Tick} (meno \emph{Delta Seconds}
        per un consumo) solo quando la luce è accesa; a zero la si spegne, e
        con \emph{Set Intensity} la si può far calare gradualmente.
  \item \textbf{Ritorno al punto sicuro.} Un vettore \emph{LastSafeLocation}
        nel personaggio, aggiornato periodicamente con \emph{Get Actor
        Location} quando \emph{Is Moving On Ground} è vero; una trigger box
        sotto il livello che, all'overlap con il personaggio, fa \emph{Set Actor
        Location} su quel vettore.
  \item \textbf{L'oggetto del 35\%.} Sul \emph{Begin Play} dell'oggetto:
        \emph{Random Bool with Weight} con peso $0{,}35$; se esce falso,
        \emph{Destroy Actor}. Farlo al Begin Play, prima che il giocatore
        possa vederlo, è ciò che rende la sparizione invisibile; per lo stesso
        motivo l'oggetto non deve lasciare tracce (collisioni, luci, suoni
        legati ad altri Actor).
  \item \textbf{Torretta.} Un \emph{Set Timer by Event} in loop che fa
        \emph{Spawn Actor} di un proiettile (un Actor con \emph{Projectile
        Movement Component}) nella direzione del giocatore; all'\emph{Event
        Hit} del proiettile, cast al personaggio e \emph{Open Level} con
        \emph{Get Current Level Name} (sezione~\ref{gd12:sec:ricarica}).
  \item \textbf{Tre spawn, poi 5 secondi.} Un intero \emph{SpawnCount}: a ogni
        pressione, se è minore di 3, \emph{Spawn Actor} e $+1$; quando arriva a
        3, un \emph{Delay} di 5 secondi lo riporta a 0. (Se la feature va
        intesa come ``un tasto genera tre oggetti di fila, poi 5 secondi di
        attesa'', un \emph{For Loop} da 0 a 2 e un booleano
        \emph{CanSpawn} risolvono la variante.)
  \item \textbf{Zona da 10 HP al secondo.} La \emph{Damage Area} e
        l'\emph{Health Component} del capitolo~\ref{gd13:cap:git-componenti}
        (sezione~\ref{gd13:sec:componenti}) sono già quasi la soluzione: al
        \emph{Begin Overlap} parte un \emph{Set Timer by Event} in loop ogni
        secondo che toglie 10; all'\emph{End Overlap} lo si ferma con
        \emph{Clear and Invalidate Timer by Handle}. In alternativa, sull'Event
        Tick, $10 \times$ \emph{Delta Seconds} finché si è dentro.
  \item \textbf{Modello a caso ogni secondo.} Un array di \emph{Static Mesh}
        (sezione~\ref{gd12:sec:array}); un \emph{Set Timer by Event} in loop da 1
        secondo che sceglie un elemento con \emph{Random} (o con \emph{Random
        Integer in Range} sull'indice) e lo passa a \emph{Set Static Mesh}.
\end{enumerate}

% ---------------------------------------------------------------------
\section{Suggerimenti durante il lavoro}
\label{gd16:sec:suggerimenti}

Mentre la classe lavora, il docente gira tra i banchi e tra gli schermi
condivisi. Tre indicazioni vengono date ad alta voce a tutti.

\subsection{Get Player Pawn}

``Potrebbe servire una cosa, quindi ve la butto lì.'' Finora la reference al
giocatore la si è sempre ottenuta dal \textbf{Begin Overlap}, quando il player
tocca una box. Ma si può fare anche \textbf{Get Player Pawn}, che ``vi dà il player
a prescindere, quindi senza che tocchi una box'': ``ci avete le informazioni del
player qui dentro''. Utile per esempio per la feature~1 (la distanza dal player)
o per la torretta, che deve sapere dove mirare senza aspettare un contatto.

\paragraph{Approfondimento: Pawn, Character e cast.} \emph{Get Player Pawn}
restituisce un \emph{Pawn} generico (con un indice, \emph{Player Index}, che in
un gioco a giocatore singolo è 0). Per la posizione basta così; per le variabili
o gli eventi definiti nel proprio personaggio (come \emph{CoinGrabbed}) serve
comunque un cast al personaggio, o il nodo gemello \emph{Get Player Character}
seguito dal cast. Romero e Sewell ricordano che un Pawn è ``il corpo fisico'' del
personaggio, mentre il \emph{Controller} che lo possiede ne è ``una specie di
cervello''.

\subsection{Il personaggio con due corpi}

A chi usa il template in prima persona e nota, nel Blueprint del personaggio,
``tipo due personaggi compenetrati'', il docente spiega che non è un errore ma
una scelta di chi l'ha fatto ``super bene'':
\begin{itemize}
  \item una mesh è quella \textbf{che vedi tu in prima persona} (braccia e arma
        davanti alla telecamera);
  \item l'altra è quella \textbf{che vedono gli altri}, in terza persona.
\end{itemize}
``Capita su molti giochi, tipo su \textbf{Overwatch}'': l'animazione che vede il
giocatore è diversa da quella che vedono gli altri, ``tendenzialmente''. ``Quindi
hanno fatto la \textbf{doppia mesh} per questo motivo.'' È anche la spiegazione
del ``personaggio un po' clippato'' notato a inizio lezione
(sezione~\ref{gd16:sec:template}).

\subsection{I tasti solo su ciò che si possiede}

L'ultimo suggerimento è una differenza importante rispetto a \textbf{Unity}. Se si
mette un evento di input, ``che ne so, il tasto qua'', \emph{sulla moneta}, non
funzionerà: ``voi non state possedendo la coin, perché voi state possedendo il
First Person Character''.
\begin{quote}
Quindi mi raccomando: i tasti possono essere messi soltanto sulla roba che
possedete, altrimenti non funziona. Su Unity invece funzionerebbe, questa roba
qua.
\end{quote}
Per le feature dell'esercitazione, quindi, gli input (jetpack, torcia, spawn)
vanno nel personaggio, che poi può chiamare eventi su altri Actor, esattamente
come la moneta chiama \emph{CoinGrabbed} ma al contrario.

\paragraph{Approfondimento: possedere un Pawn.} Romero e Sewell (capitolo~3)
spiegano il meccanismo: le istanze di Pawn e Character ricevono gli eventi di
input \textbf{solo se un PlayerController le possiede}, e gli eventi di input
si possono mettere nel PlayerController o nel Pawn (il libro consiglia di
scegliere una strada e di restare coerenti nel progetto). Con il nodo
\emph{Possess} si può cambiare a runtime il Pawn controllato, e da quel momento
solo il nuovo riceve i comandi. Per completezza: un Actor qualunque può
ricevere input se glielo si chiede esplicitamente, con il nodo \emph{Enable
Input} o con l'opzione \emph{Auto Receive Input} nei suoi dettagli; ma è
un'eccezione da usare con giudizio, e la regola del docente resta la più sicura.

% ---------------------------------------------------------------------
\section{La chiusura}
\label{gd16:sec:chiusura}

La lezione prosegue con il lavoro sulle feature fino alla fine, e la
registrazione non ne restituisce i dettagli. L'ultima frase leggibile è rivolta
a chi segue da remoto: ``per quelli online, che si recupereranno la lezione la
prossima volta\dots''\footnote{Nella trascrizione ``Tobiquelli online che si
recupereranno la lezione la prossima volta, anche se'': la frase è tagliata e
la registrazione finisce lì.}. Il prossimo appuntamento con la programmazione
porterà il task con il \textbf{best score} da salvare, annunciato a metà
lezione.

% ---------------------------------------------------------------------
\section*{Punti chiave del capitolo}
\addcontentsline{toc}{section}{Punti chiave del capitolo}

\begin{itemize}
  \item Unreal Engine~6 è atteso a breve; le versioni escono circa ogni sei
        mesi. La build per browser, tolta anni fa e rimessa nella 5.9, non
        poteva essere un plugin: è una limitazione dell'engine.
  \item Secondo il docente, con l'IA capace di scrivere un engine intero torna
        possibile farsi \textbf{l'engine per il proprio gioco}, senza il peso
        (\emph{bloat}) e i costi di quelli generalisti; Godot oggi è debole nel
        3D.
  \item Il template \textbf{in prima persona} è come quello in terza; se manca,
        lo si installa dal launcher (\emph{template e pacchetti di
        funzionalità}), che permette anche di risparmiare spazio.
  \item Non esiste un ``salva tutto'': bisogna dire all'engine \textbf{quali
        variabili} salvare e come. Che cosa salvare è anche una decisione di
        design.
  \item La classe \textbf{Save Game} (es. \emph{PlayerSaves}) è un Blueprint
        senza viewport di cui contano solo le \textbf{variabili}; è l'unico
        modo, oltre al C++, di scrivere dati permanenti sul disco. Di solito
        una o due per gioco: partita e impostazioni.
  \item La \textbf{Game Instance} (una sola, prima a caricarsi, ultima a
        scaricarsi, sopravvive ai cambi di mappa) è il posto per ciò che è
        globale: salvataggi, achievement, musica continua. Il suo \textbf{Event
        Init} è il Begin Play del gioco intero.
  \item All'avvio: \textbf{Does Save Game Exist} $\rightarrow$ \textbf{Branch};
        se no \textbf{Create Save Game Object}, se sì \textbf{Load Game from
        Slot} + \textbf{Cast To PlayerSaves}; in entrambi i casi \textbf{Set} in
        una variabile reference \emph{Save} del tipo del salvataggio.
  \item Lo \textbf{slot name} è il nome del file: sempre lo stesso, mai vuoto,
        comodo da tenere in una variabile String.
  \item Creare non è salvare: il file nasce con \textbf{Save Game to Slot},
        chiamato da un custom event \textbf{SaveGame} nella Game Instance,
        raggiungibile da ovunque. Le varianti \textbf{async} non bloccano il
        gioco (l'icona che lampeggia) e segnalano la fine con
        \emph{Completed}.
  \item Ogni Blueprint è anche un \textbf{tipo di variabile}: per raggiungere lo
        score si passa per Get Game Instance $\rightarrow$ cast $\rightarrow$
        Save $\rightarrow$ Score, una ``torre'' di referenze; il cast si può
        rendere \textbf{puro}.
  \item La logica sta in chi ha senso che la gestisca: la moneta avvisa il
        personaggio con un \textbf{custom event} (\emph{CoinGrabbed}) e si
        distrugge; il personaggio aggiorna il salvataggio. I custom event
        servono a tenere ordinato il codice.
  \item Il file è in \texttt{Saved/SaveGames/Save1.sav} nel progetto (nella
        build, in una cartella dell'utente). Posizioni e oggetti complessi si
        salvano con Vector o Transform e codice dedicato al ripristino.
  \item Per vedere un valore sempre a schermo senza widget: Print String
        sull'\textbf{Event Tick} con durata zero.
  \item \textbf{Get Player Pawn} dà il giocatore senza bisogno di un overlap. Il
        personaggio in prima persona ha \textbf{due mesh} (quella vista da sé e
        quella vista dagli altri). Gli \textbf{input} funzionano solo sui Pawn
        posseduti, a differenza di Unity.
  \item Nel prossimo task si salverà il \textbf{best score}.
\end{itemize}
